\documentclass[10pt,a4paper,onecolumn]{article}
\usepackage{course-paper}
\newcommand{\papertitle}{Image Restoration by Deconvolution}
\newcommand{\paperauthors}{Pham Manh Dung, Nguyen Minh Hieu, Vu Dinh Hieu, and Nguyen Thanh Tam}
\newcommand{\institution}{Posts and Telecommunications Institute of Technology}
\newcommand{\classcode}{E23TTNT02}
\newcommand{\instructor}{Nguyen Tat Thang}

\newcommand{\RunID}{bsds500\_20261118\_234602}
\newcommand{\PythonVersion}{3.13.5}
\newcommand{\NumPyVersion}{2.5.4}
\newcommand{\SciPyVersion}{1.16.0}
\newcommand{\JVMVersion}{25.0.1}
\newcommand{\TestIDs}{100007, 100039, 100099, 10081, 101027, 101084, 102062, 103006, 103029, 103078}
\newcommand{\ValidationIDs}{101085, 101087, 102061, 103070, 105025}
\newcommand{\ProbeError}{$2.90\times 10^{-8}$}
\newcommand{\PythonMetricError}{$4.35\times 10^{-8}$}
\newcommand{\MetricPSNRGap}{$3.03\times 10^{-12}$}
\newcommand{\MetricSSIMGap}{$4.63\times 10^{-14}$}
\newcommand{\MetricMSEGap}{$2.57\times 10^{-14}$}
\newcommand{\NativePSNRGap}{$2.84\times 10^{-13}$}
\newcommand{\BenchmarkSpeedRange}{1.77--3.01}
\newcommand{\BenchmarkMaxError}{$2.89\times 10^{-15}$}
\newcommand{\EpflProjectionError}{$8.58\times 10^{-8}$}
\newcommand{\EpflMassPercent}{16.55}
\newcommand{\EpflParityError}{$5.13\times 10^{-6}$}

\hypersetup{pdftitle={\papertitle},pdfauthor={\paperauthors}}

\begin{document}
\hypersetup{pageanchor=false}
\newgeometry{left=30mm,right=30mm,top=33mm,bottom=29mm}
\begin{titlepage}
  \AddToShipoutPictureBG*{\AtPageLowerLeft{%
    \hspace*{12mm}\raisebox{12mm}{\includegraphics[
      width=\dimexpr\paperwidth-24mm\relax,height=\dimexpr\paperheight-24mm\relax]{background.png}}}}
  \centering
  {\large\bfseries \MakeUppercase{\institution}\par}
  \vspace{.9cm}
  \includegraphics[width=35mm]{logo.png}\par
  \vspace{1.05cm}
  {\large\scshape Course Project Report\par}
  \vspace{.2cm}
  {\large Digital Image Processing\par}
  \vspace{1.0cm}
  {\Huge\bfseries \papertitle\par}
  \vspace{.4cm}
  {\Large Implementation in Python and Validation\\[.15em]
  with DeconvolutionLab2 in ImageJ\par}
  \vspace{.55cm}
  \rule{24mm}{.6pt}\par
  \vfill
  {\large
  \renewcommand{\arraystretch}{1.18}
  \begin{tabular}{@{}ll@{}}
    \textbf{Instructor} & \instructor\\
    \textbf{Class} & \classcode
  \end{tabular}\par
  \vspace{.6cm}
  \begin{tabular}{@{}l@{\hspace{1.35cm}}l@{}}
    \toprule\textbf{Student ID} & \textbf{Student}\\\midrule
    B23DCCE024 & Pham Manh Dung\\
    B23DCCE030 & Nguyen Minh Hieu\\
    B23DCCE036 & Vu Dinh Hieu\\
    B23DCCE084 & Nguyen Thanh Tam\\
    \bottomrule
  \end{tabular}\par}
  \vfill
  {\large November 2026\par}
\end{titlepage}
\restoregeometry

\pagenumbering{arabic}
\hypersetup{pageanchor=true}

\begin{flushleft}
{\LARGE\scshape \papertitle\par}
\vspace{.68em}
{\normalsize\bfseries \paperauthors\par}
\vspace{.25em}
{\small \institution\par
Digital Image Processing \enspace\textbar\enspace Class \classcode\par}
\end{flushleft}

\pdfbookmark[1]{Abstract}{abstract}
\begin{center}\large\scshape Abstract\end{center}
\begin{quote}
Non-blind image deconvolution recovers an image from blurred, noisy measurements when the point spread function (PSF) is specified. We derive stabilized inverse filtering, two Tikhonov estimators, Richardson--Lucy (RL), and its total-variation variant (RLTV) from a common forward model. A Python implementation is tested against independent numerical oracles and DeconvolutionLab2 2.1.2 in ImageJ. The main experiment uses five BSDS500 validation images to select parameters and ten disjoint test images under Gaussian and Poisson noise. On matched Gaussian blur, Python RLTV achieves mean PSNRs of \(25.44\) and \(24.89\) dB, gains of \(1.76\) and \(1.92\) dB over the observations. After harmonizing conventions, four solvers agree with the plugin within \(2.3\times10^{-5}\) intensity units on these cases. Sensitivity experiments vary PSF width, blur type, and boundaries. An independent EPFL Hollow Bars projection tests a much wider optical PSF and the transfer of parameters selected on natural images. Raw TIFFs, independent MSE/PSNR/SSIM measurements, and per-image results accompany the study.
\end{quote}
\textbf{Keywords}\quad image restoration, point spread function, Fourier transform, Tikhonov regularization, Richardson--Lucy, ImageJ

\textbf{Code availability}\quad The \href{https://github.com/Hieuvu4438/image-deconvolution-btl/tree/main/code}{\textcolor{listingkeyword}{\underline{GitHub source code folder}}} contains the implementations, experiment scripts, and numerical tests.

\section{Introduction}
\label{sec:intro}

The image recorded by an acquisition system is an indirect measurement of a scene. Optical blur, motion, defocus, and sampling can spread the intensity of a point across nearby pixels. Noise perturbs the resulting measurements. Even with a known blur kernel, restoring attenuated spatial frequencies can also amplify noise. A method that sharpens a contour can therefore create ringing, increase error on a smooth background, or yield intensities outside the intended range.

This project studies non-blind deconvolution of single-channel images. It develops the image formation model and implements five estimators in Python. DeconvolutionLab2 provides an independently implemented reference in ImageJ. For numerical comparison, the PSF origin, boundary operator, FFT precision, normalization, initialization, and clipping matter. Measuring quality independently on the same saved pixels helps distinguish a difference in reconstruction from a difference in scoring.

\subsection{Related work}

\citet{tikhonov} established regularization for ill-posed problems. \citet{hansen} show how boundary conditions determine the structure of the blur matrix and how this structure enables spectral filtering. \citet{fessler} discusses the trade-off between data fidelity and image priors in regularized reconstruction. We use identity and Laplacian penalties to illustrate the role of the regularization operator.

\citet{richardson} and \citet{lucy} developed multiplicative iterative methods for recovering nonnegative distributions. The Poisson formulation also follows from an expectation--maximization (EM) allocation of unobserved counts, as in the linear model of \citet{shepp}. \citet{dey} combined Richardson--Lucy with total variation to control noise in confocal deconvolution. Likelihood in the measurement domain and reconstruction accuracy in the image domain can favor different iteration counts.

\citet{sage} present DeconvolutionLab2 and evaluate its algorithms on synthetic and acquired microscopy data. Their work motivates measuring numerical quality, visible artifacts, and computational cost. We use BSDS500 \citep{bsds} for natural image structure and the public EPFL Hollow Bars volume \citep{epflbars} for an independent structured object and a theoretical microscope PSF.

\subsection{Research questions}

\begin{enumerate}
\item \textbf{Reconstruction quality.} Under known blur, how do direct and iterative estimators change full-reference error and image structure? How does performance depend on the regularization parameter or iteration limit?
\item \textbf{Numerical agreement.} When Python and DeconvolutionLab2 solve the same discrete problem, how close are their raw outputs? Which named counterparts retain different discrete rules?
\item \textbf{Robustness.} What happens when PSF width, blur type, or the boundary extension is wrong, or when parameters selected on natural photographs are transferred to a much broader optical blur?
\end{enumerate}

\subsection{Scope and contributions}

This project contributes a reproducible implementation and a controlled evaluation of established algorithms. It separates calibration and test images, uses a shared float32 handoff, computes quality independently in Java and Python, and reports raw and clipped scores. Complementary experiments on a geometric phantom and an axial-sum projection from EPFL probe structures that a global statistic can obscure.

All reconstructions are two-dimensional and assume a specified, spatially invariant blur. The EPFL volume is synthetic, and its projection supplies a controlled 2D benchmark. Appendix~\ref{app:reproduce} gives the reproduction procedure and Appendix~\ref{app:atlas} documents every main test image under both noise models.

\section{Forward model and Fourier-domain formulation}
\label{sec:model}

\subsection{From a point response to a discrete operator}

Let \(f\in\mathbb R^{M\times N}\) denote the sharp intensity image and \(g\) the observation on the same pixel grid. In a linear, shift-invariant system the response to a unit source at position \((p,q)\) is a shifted point spread function (PSF). Expanding the image as a weighted sum of impulses therefore gives a weighted sum of shifted PSFs. This is convolution. For an integer PSF support and periodic extension, the forward map is
\begin{equation}
(\Hoperator f)[i,j] = \sum_{a,b} h[a,b]\,
    f[(i-a)\bmod M,(j-b)\bmod N].
\label{eq:conv}
\end{equation}
Indices \(a,b\) are offsets from the PSF origin. The intensity PSFs in this work are nonnegative and sum to one. This convention keeps constant images fixed and preserves total intensity before adding noise. The two conservation statements follow from the normalization and the use of all periodic shifts. The same PSF under a cropped zero-boundary operator can lose intensity at the edges.

For a stored kernel with shape \(m\times n\), its origin is at \((\lfloor m/2\rfloor,\lfloor n/2\rfloor)\) in zero-based indexing. This definition is explicit for even as well as odd sizes. To construct the image-sized impulse response, place the stored kernel in the upper-left block of a zero array and roll by the negative origin in both dimensions. Padding a centered kernel and shifting the whole image-sized array also works when the center is placed exactly. A one-pixel placement error changes the phase, even if a magnitude spectrum appears unchanged.

For a \(1\times 5\) signal \((0,0,8,0,0)\), the symmetric kernel \(h[-1]=1/4,\ h[0]=1/2,\ h[1]=1/4\) gives \((0,2,4,2,0)\). For a \(5\times5\) image containing a single intensity-9 point, the uniform \(3\times3\) kernel gives a \(3\times3\) block of ones around that point. Moving the point to a corner sends part of the block to the opposite boundaries under periodic extension. These examples motivate the impulse checks in Appendix~\ref{app:oracles}.

\subsection{Boundary conditions and the adjoint}

Periodic extension identifies opposite edges. Zero extension assumes that everything beyond the field of view has value zero. Reflection extends the last visible structures symmetrically. Each assumption changes which pixels contribute to a boundary measurement and hence changes the inverse problem. An FFT on the original grid computes circular convolution. To compute full linear convolution with a finite kernel by FFT, pad to at least \((M+m-1)\times(N+n-1)\), then crop the intended observation and account for the kernel origin. Padding for a forward computation alone does not define the adjoint of the resulting cropped operator.

The adjoint \(\Hoperator^\top\) is defined by
\(\langle \Hoperator f,u\rangle=\langle f,\Hoperator^\top u\rangle\).
For circular convolution, changing \(i-a\) and \(j-b\) to summation coordinates in the inner product gives circular convolution with \(h[-a,-b]\). For a real centered symmetric kernel, the forward and adjoint maps coincide. An asymmetric impulse response makes a stronger test because it distinguishes convolution from correlation and exposes a sign error in the shift. With a windowed measurement operator, the adjoint must also reverse the crop by embedding the measurement into the larger grid.

\subsection{PSFs used by the experiments}

For an isotropic Gaussian, the finite samples on \([-8,8]^2\) are proportional to
\(\exp[-(a^2+b^2)/(2\sigma_h^2)]\). The canonical width is \(\sigma_h=1.9\) pixels. The full width at half maximum of the corresponding continuous profile is \(2\sqrt{2\log 2}\sigma_h\), approximately \(4.47\) pixels. Normalization is applied after sampling and truncation. The horizontal motion PSF places mass \(1/9\) on each displacement \(b=-4,\ldots,4\) at \(a=0\). It represents uniform translation during exposure.

An actual optical PSF can be estimated by imaging fluorescent beads below the system resolution or predicted from an optical model. Numerical aperture, wavelength, refractive index, focus, and physical sampling all affect the kernel. The EPFL PSF Generator includes the Born--Wolf and Gibson--Lanni models \citep{psfgen,sage}. The independent Hollow Bars experiment uses a supplied widefield PSF rather than a Gaussian approximation. The model assumes that the kernel has already been calibrated on the imaging grid. Estimating a spatially varying PSF is outside its scope.

\subsection{Observation noise}

Write \(b=\Hoperator f\). An additive model is \(g=b+n\), often with independent zero-mean Gaussian \(n_i\) and common variance \(\sigma_n^2\). Unweighted squared residuals correspond to the negative Gaussian log-likelihood up to a scale and an additive constant. In our simulation the Gaussian observation is clipped from below at zero so the exact likelihood of the clipped measurement differs near the dark end.

For independent photon counts, \(y_i\mid f\sim\operatorname{Pois}(P b_i)\), where \(P\) denotes the count scale at normalized intensity one. Dividing by \(P\) gives \(g_i=y_i/P\) with
\begin{equation}
\mathbb E[g_i\mid f]=b_i,\qquad
\operatorname{Var}(g_i\mid f)=b_i/P .
\end{equation}
The variance changes with the signal. A sample may exceed one in either simulation. The generation stage does not cap the bright end. If a detector adds known background \(r_i\) in count units, the Poisson mean becomes \(P b_i+r_i\). A correct likelihood would include \(r_i\) inside its logarithm. Our controlled Poisson data do not include a background term.

\subsection{The optical transfer function and conditioning}

Use a capital letter for a 2D discrete Fourier transform (DFT). The optical transfer function (OTF) is \(H=\F h\), with the origin at array index \((0,0)\). The convolution theorem gives
\begin{equation}
G[u,v]=H[u,v]F[u,v]+\mathcal N[u,v],\qquad H[0,0]=1
\label{eq:dft}
\end{equation}
under the additive model. The adjoint uses \(\overline H\). Padding with a different origin changes phase through the shift theorem, while leaving \(|H|\) unchanged.

Let \(U\) be the unitary DFT on vectorized images. The circular blur matrix is block circulant with circulant blocks and has decomposition
\(\Hoperator=U^*\operatorname{diag}(H)U\) \citep{hansen}. Its singular values are \(|H[u,v]|\). When all are nonzero, the spectral condition number is \(\max|H|/\min|H|\). Normalized nonnegative PSFs have \(\max|H|=1\) by the triangle inequality and the DC value. Thus the smallest transmitted magnitude controls the worst-case noise amplification. If \(H=0.02\), a spectral perturbation of \(0.01\) becomes an inverse error of \(0.5\).

If \(H=0\) at a bin, any component of \(F\) in that direction leaves the data unchanged. A two-pixel average maps both \((1,0)\) and \((0,1)\) to \((1/2,1/2)\); the alternating component is lost. The Moore--Penrose solution sets null-space coefficients to zero. The box-motion OTF has regularly spaced zeros as a continuous function on the frequency torus. A finite DFT may sample a zero or pass close to it, depending on the grid size. Neither stronger sharpening nor more iterations supplies information at a true null.

\begin{figure}[!htbp]
\centering\includegraphics[width=\linewidth]{frequency.png}
\caption{Horizontal OTF slices for the two simulation PSFs, and filter gains on the Gaussian slice. The inverse gain is clipped only for plotting the logarithmic scale.}
\label{fig:frequency}
\end{figure}

\section{Restoration algorithms}
\label{sec:algorithms}

\subsection{Stabilized inverse filtering}

At nonzero \(H\), the algebraic inverse is \(\widehat F=G/H=F+\mathcal N/H\). We use a stabilized variant called NIF, matching the usage in this study and the plugin:
\begin{equation}
\widehat F_{\mathrm{NIF}} =
\frac{\overline H\,G}{\max(|H|^2,\tau)},\qquad \tau=10^{-4}.
\label{eq:nif}
\end{equation}
For \(|H|\ge\sqrt{\tau}\), the response is the exact inverse. Below that threshold its gain is \(|H|/\tau\), which tends to zero as \(|H|\) approaches zero. The resulting maximum gain is \(1/\sqrt{\tau}=100\). Dropping frequency bins entirely would be a different stabilization rule. This NIF is a deliberate baseline for measuring how far finite inverse gain alone protects against noise. It can yield negative or saturated intensities because it does not use spatial constraints.

\subsection{Tikhonov estimators}

The Tikhonov family balances a data mismatch and a quadratic penalty:
\begin{equation}
\widehat f_\lambda=
\arg\min_f\left(
\tfrac12\norm{\Hoperator f-g}_2^2+\tfrac{\lambda}{2}\norm{Lf}_2^2
\right),\quad \lambda>0 .
\label{eq:tik-objective}
\end{equation}
The gradient is
\(\Hoperator^\top(\Hoperator f-g)+\lambda L^\top Lf\).
Setting it to zero yields the normal equation
\begin{equation}
(\Hoperator^\top\Hoperator+\lambda L^\top L)\widehat f_\lambda
=\Hoperator^\top g .
\end{equation}
The coefficient matrix is positive definite exactly when
\(\ker\Hoperator\cap\ker L=\{0\}\), since its quadratic form is a sum of two squared norms. If \(L\) is periodic and has DFT response \(D\), both terms diagonalize, giving
\begin{equation}
\widehat F_\lambda=\frac{\overline H G}{|H|^2+\lambda |D|^2}.
\label{eq:tikhonov}
\end{equation}
Equivalently, under a Gaussian measurement model this is a quadratic-prior estimate with the noise and prior scales absorbed into \(\lambda\). A penalty operator with a null space leaves those directions unpenalized; uniqueness still holds if the data identify them.

\subsubsection{Identity penalty (TRIF)}

With \(L=I\), every frequency has penalty \(\lambda\). Its DC response is \(1/(1+\lambda)\), so the method reduces the mean even if the blur preserves it. For \(H=0.1\), \(G=0.45\), and \(\lambda=0.01\), the inverse gives \(4.5\), while TRIF gives \(0.045/0.02=2.25\). The correction of an unstable bin is thus paid for in bias. The response to noise is \(x/(x^2+\lambda)\) for \(x=|H|\), bounded by \(1/(2\sqrt{\lambda})\) by the arithmetic--geometric mean inequality. The bound at \(\lambda=0.01\) is 5, compared with 100 for the NIF setting.

\subsubsection{Laplacian penalty (RIF)}

Use the five-point stencil with center \(-4\) and four unit neighbors. Its spectral response is
\begin{equation}
D[u,v]=-4\left(\sin^2(\pi u/M)+\sin^2(\pi v/N)\right).
\label{eq:laplacian}
\end{equation}
DC has \(D=0\) and remains unchanged because \(H[0,0]=1\). At higher frequencies the penalty usually grows, discouraging oscillatory estimates. DeconvolutionLab2 uses a spectral Laplacian \(\omega_x^2+\omega_y^2\) on the even grids considered here. The two Laplacians agree in leading order near DC, up to sign, but their magnitudes differ at higher frequencies. Using the same numerical \(\lambda\) therefore need not implement the same restoration.

\subsection{Bias, variance, mismatch, and parameter choice}

For a correct forward model, define \(W_\lambda=\overline H/(|H|^2+\lambda |D|^2)\) and \(A_\lambda=W_\lambda H\). Then
\begin{equation}
\widehat F_\lambda-F=(A_\lambda-1)F+W_\lambda\mathcal N .
\label{eq:bias}
\end{equation}
With zero-mean uncorrelated additive noise and given \(F\), expected squared error in one bin is
\(|A_\lambda-1|^2|F|^2+|W_\lambda|^2 S_n\), where \(S_n\) is noise power. Increasing \(\lambda\) weakens gains and can remove real detail. Under stationary signal and noise models the Wiener response has the same algebraic form when \(\lambda |D|^2=S_n/S_f\) for positive signal power \(S_f\) \citep{hansen,sage}. Here the penalty is chosen by its structure and the scalar parameter by validation.

If the filter \(W\) was designed using \(\widetilde H\) but the true system is \(H\), the residual error includes \(W(H-\widetilde H)F\). Consequently PSF errors can produce ringing or shape distortion even with zero measurement noise. In applications with no reference image, a discrepancy principle can set the data residual to match the expected measurement noise scale when that scale and the forward operator are credible \citep{fessler}. Our protocol uses held-out reference images to select a practical common configuration. It does not fit to the test set.

\subsection{Richardson--Lucy from a Poisson likelihood}

The negative Poisson log-likelihood divided by \(P\), ignoring terms independent of the estimate, is
\begin{equation}
\Phi(f)=\sum_i\left[(\Hoperator f)_i-g_i\log((\Hoperator f)_i)\right],
\quad f\ge0.
\end{equation}
The gradient where the predicted intensities are positive is
\(\Hoperator^\top(\mathbf 1-g/(\Hoperator f))\). To derive the multiplicative update, let \(A_{ij}\) represent the PSF contribution from true pixel \(j\) to measured pixel \(i\). The latent count \(z_{ij}\) has mean \(P A_{ij}f_j\), and \(y_i=\sum_j z_{ij}\) is observed. Conditional allocation of independent Poisson variables and the EM maximization step give
\begin{align}
\mathbb E[z_{ij}\mid y_i,f^{(k)}] &=
y_i\frac{A_{ij}f^{(k)}_j}{(\Hoperator f^{(k)})_i},\\
f^{(k+1)}_j &=
\frac{\sum_i\mathbb E[z_{ij}\mid y_i,f^{(k)}]}
{P\sum_i A_{ij}} .
\end{align}
Using \(g_i=y_i/P\) yields \citep{richardson,lucy,shepp}
\begin{equation}
f^{(k+1)}= f^{(k)}\odot
\frac{\Hoperator^\top(g/(\Hoperator f^{(k)}))}{\Hoperator^\top\mathbf 1}.
\label{eq:rl}
\end{equation}
The division and product are pointwise. The denominator is one for a circular unit-sum PSF. With valid positive predictions at nonzero measurements, exact arithmetic gives
\(\sum_j f_j^{(k+1)}=\sum_i g_i\)
after collecting terms by observation pixel. Plain RL preserves nonnegativity for nonnegative inputs. A zero in the initialization remains zero under the multiplicative rule, which matters for images initialized from a clipped observation.

For the two-pixel blur with diagonal \(3/4\) and off-diagonal \(1/4\), let \(g=(6,2)^\top\), \(f^{(0)}=(4,4)^\top\). The first prediction is \((4,4)^\top\) and the ratio is \((3/2,1/2)^\top\). Applying the adjoint and multiplication gives \(f^{(1)}=(5,3)^\top\). A second step gives \((40/7,16/7)^\top\). The exact inverse of \(g\) is \((8,0)^\top\), while increasing the first observation by 0.1 and reducing the second by 0.1 gives the inverse \((8.2,-0.2)^\top\). Appendix~\ref{app:oracles} records the independent checks used for such examples.

In the exact EM setting the data likelihood is nondecreasing. This concerns predictions in the observation domain and does not imply an increasing PSNR against the true image. Iterations may progressively reproduce fine features, then fit statistical fluctuations. We therefore treat \(K\) as an implicit regularization parameter. The Python main mode starts from the observed mean and guards its prediction with a nonnegative floor and an additive \(10^{-12}\). The plugin starts from \(g\) and returns zero ratio below its epsilon threshold. The auxiliary Python matching mode reproduces those rules. Such guards are numerical choices and can change exact conservation and likelihood identities.

\subsection{Total-variation regularized Richardson--Lucy}

We consider the penalized objective \(\Phi(f)+\lambda\,\operatorname{TV}_\delta(f)\), following the use of TV in Poisson microscopy \citep{dey,sage}. With the forward periodic gradient \(D_+\),
\begin{equation}
\operatorname{TV}_\delta(f)=\sum_i
\sqrt{\norm{(D_+f)_i}_2^2+\delta^2}.
\end{equation}
The two gradient components share one isotropic magnitude. The smoothing constant \(\delta=0.01\) prevents division by zero on flat regions. Define \(\operatorname{div}=-D_+^\top\). Differentiating gives
\(-\operatorname{div}(D_+f/\sqrt{\norm{D_+f}^2+\delta^2})\),
where normalization uses the magnitude at each pixel. For the present periodic system, the implemented iteration uses
\begin{align}
v^{(k)}_i &= \frac{(D_+ f^{(k)})_i}
{\sqrt{\norm{(D_+ f^{(k)})_i}_2^2+\delta^2}},\\
f^{(k+1)} &= f^{(k)}\odot
\frac{\Hoperator^\top\big(g/(\Hoperator f^{(k)}+\varepsilon)\big)}
{1-\lambda\operatorname{div}v^{(k)}} .
\label{eq:rltv}
\end{align}
Under this discretization \(|\operatorname{div}v_i|\le4\), so \(\lambda<1/4\) is sufficient for positive denominators. The code verifies positivity and checks for nonfinite iterates. Setting the penalty to zero reduces the scheme to RL under the same numerical guards. TV can preserve step boundaries while suppressing oscillatory background, but excessive weight can erase fine textures and cause piecewise-flat artifacts. The finite-step rule with stabilization is evaluated empirically. We do not assert that every step minimizes the penalized objective.

\begin{table}[!htbp]
\centering
\caption{The controls and principal failure modes of the five implemented estimators.}
\small
\begin{tabularx}{\linewidth}{@{}lYY@{}}
\toprule Method & Control & Failure mode to inspect\\\midrule
NIF & Power floor \(\tau\) & Strong noise amplification in weakly transmitted bands.\\
TRIF & Identity penalty weight \(\lambda\) & Intensity bias and loss of structure if too large.\\
RIF & Laplacian and weight \(\lambda\) & Oscillatory edges with weak regularization or a wrong PSF.\\
RL & Iteration limit and initialization & Slow recovery, zero locking, or increased noise after many iterations.\\
RLTV & \(K,\lambda,\delta\) and discrete TV & Texture loss or different results under a different TV operator.\\
\bottomrule
\end{tabularx}
\end{table}

\section{Implementation and numerical verification}
\label{sec:implementation}

\subsection{Python architecture}

The entry point \file{solve} accepts a finite 2D observation, a centered unit-sum PSF, and a method identifier. It returns a floating-point estimate on the same grid. The input contract catches shape mismatch, unsupported algorithm names, nonfinite values, invalid parameter ranges, negative PSF coefficients, and negative observations in the multiplicative solvers. Conversion to a display range happens at scoring or visualization. It is separate from the reconstruction itself.

\begin{table}[!htbp]
\centering
\caption{Implementation components and their normal entry points relative to the project root.}
\small
\begin{tabularx}{\linewidth}{@{}p{.38\linewidth}Y@{}}
\toprule Component & Responsibility\\\midrule
\file{code/deconv/core.py} & TIFF/PNG handoff, PSF construction, circular blur, and measurement simulation.\\
\file{code/deconv/solvers.py} & Direct spectral filters, iterative estimators, and the RL rule used for matching the plugin.\\
\file{code/deconv/metrics.py} & MSE, fixed-peak PSNR, and weighted local SSIM.\\
\file{code/scripts/} & Dataset preparation, experiment orchestration, audit, and report inputs.\\
\file{code/imagej/BatchReference.java} & ImageJ TIFF reading, plugin execution, independent metric calculation, and CSV export.\\
\file{code/tests/test_deconv.py} & Analytical examples and independent numerical oracles.\\
\bottomrule
\end{tabularx}
\end{table}

Listing~\ref{lst:forward-code} shows how the forward operator places the PSF origin at index \((0,0)\) before computing the OTF. The inverse \file{irfft2} receives the original image shape so it also reconstructs an odd last dimension correctly. The solver uses the same OTF for blurring and its complex conjugate for back-projection. Listing~\ref{lst:spectral-code} computes the Laplacian spectrum directly from the row and column frequencies, including for one-row and one-column inputs.

\reportlisting{Python}{deconv/core.py}{46}{63}{
  caption={PSF centering, validation, and circular convolution in \texttt{core.py}.},
  label={lst:forward-code}
}

\subsection{Iterative update and stability checks}

Listing~\ref{lst:iterative-code} shows the RL and RLTV loop. The main branch clamps negative roundoff in the prediction and correction and uses \(\varepsilon=10^{-12}\). The plain RL matching branch uses the plugin's observation initialization and zero-ratio threshold. RLTV checks the denominator before division and verifies that every updated value is finite. Appendix~\ref{app:code} includes the full \file{solve} implementation and the experiment and scoring excerpts.

For \(Q=MN\), the direct methods take \(O(Q\log Q)\) time and \(O(Q)\) working memory. RL and RLTV take \(O(KQ\log Q)\) time, with \(O(Q)\) extra work per iteration for the TV terms. The constant can be substantial because the prediction and correction each require a forward and inverse transform. Reusing the PSF spectrum reduces unnecessary FFT calls without changing the iteration count. The code distinguishes the Python float64 arithmetic from the float32 handoff to ImageJ.

\subsection{ImageJ reference workflow}

The Java batch launcher invokes \file{DeconvolutionLab2_Run} from DeconvolutionLab2 2.1.2 on ImageJ 1.54p \citep{sage,dl2docs}. Listing~\ref{lst:imagej-code} shows the plugin invocation and scoring handoff. Appendix~\ref{app:imagej} gives the corresponding settings for replay in the GUI.

The launcher sets macro options rather than typing into a GUI. It reads exported TIFFs through ImageJ, checks their bit depth and dimensions, and independently computes the metrics. The reference runtime is Java \JVMVersion. The batch records the command for every case, the JAR URLs and SHA-256 hashes, and the algorithm outputs as float TIFF. It also calls the plugin's \file{Assessment.psnr} for an additional native-score check. The JAR and the Python source may use different single-precision normalization and discrete derivative rules, so the settings and source comparison are part of the evidence \citep{dl2code}.

\subsection{Verification strategy}

The 20 Python tests cover hand calculations, agreement with spatial convolution for an asymmetric PSF, odd/even centering, a dense Tikhonov oracle, the adjoint in RL, a directional derivative of TV, SSIM from an explicit window, reproducible simulation, raw TIFF round-trip, and equivalence of the two FFT branches. The one-row Laplacian test protects the direct spectral implementation against support-size errors. A separate one-step oracle tests the plugin's zero-ratio convention.

The main test audit checks the exact key set for 240 Python records, then reads every referenced TIFF and recomputes the metrics. The maximum metric difference due to the float32 export is \PythonMetricError. A second batch verifies the 240 main plugin metric records plus two for an asymmetric impulse check. The impulse is placed at the border so a wrong phase or missing wrap is observable. Its largest difference from the Python result is \ProbeError. For plugin scoring, both implementations read the same saved pixels. This makes the metric comparison logically independent of the restoration comparison.

Computational measurements time only \file{solve}, excluding TIFF handling, metric calculation, figures, and JVM overhead. They alternate the two FFT modes after warm-up on one validation input. A separate \file{tracemalloc} pass records tracked allocations. This design supports statements about those Python implementations on that machine, while Java serves as a numerical reference.

\section{Experimental protocol}
\label{sec:protocol}

\subsection{Natural images and ground truth}

BSDS500 contains natural photographs and human annotations for contour detection and segmentation \citep{bsds,bsdsdata}. We use the photographs as reference intensity images for a different task. The human labels play no role. We select the first five validation and first ten test JPEGs in lexical filename order, before fitting any parameters. A manifest records the archive checksum, original members, IDs, split, extracted locations, and SHA-256 for each JPEG. Appendix~\ref{app:reproduce} lists the exact identifiers.

Each image is decoded with Pillow, converted to 8-bit grayscale mode \file{L}, center-cropped to \(256\times256\), divided by 255, and represented as float32 for the shared TIFF handoff. These are encoded natural photographs rather than calibrated linear-light luminance. The known reference for the experiment is the resulting grayscale crop. This avoids comparing estimates to an image that has separately undergone interpolation or different grayscale conversion.

\subsection{Paired degradation and frozen settings}

The main blur uses the circular Gaussian PSF with support \(17\times17\), \(\sigma_h=1.9\) pixels, and total mass one. For each reference, create two observations from the same blurred image \(b\):
\begin{align}
g_i^{\mathrm G} &= \max(0,b_i+\eta_i),
\quad \eta_i\sim\mathcal N(0,(3/255)^2),\\
g_i^{\mathrm P} &= y_i/700,\quad y_i\sim\operatorname{Pois}(700b_i).
\end{align}
The Gaussian seed is \(20261118+\mathrm{ID}\). The Poisson seed adds \(10^6\). The random number generator is NumPy's default generator in the recorded environment. The observations are rounded to float32 once, then the exact values are passed to both platforms. The PSF file is also float32 and is interpreted as a normalized kernel. We save the noise-free blur as a diagnostic and use the noisy observation as the baseline when scoring improvement.

We select each method's settings separately for each noise model. The objective is mean clipped PSNR on five validation photographs. TRIF tries \(10^{-4},10^{-3},10^{-2},0.02\). RIF tries \(10^{-4},10^{-3},0.005,0.02\). RL tries 5, 10, 20, 40, and 80 iterations. RLTV evaluates all pairs of \(K\in\{10,20,40\}\) and \(\lambda\in\{10^{-4},0.001,0.003,0.01\}\), with \(\delta=0.01\). NIF uses its fixed floor. We store the selected configuration before producing test reconstructions.

\begin{table}[!htbp]
\centering
\caption{Parameters selected on the validation split and frozen for the main test.}
\label{tab:settings}
\begin{tabular}{lcc}
\toprule Method & Gaussian noise & Poisson noise\\\midrule
NIF & $\tau=10^{-4}$ & $\tau=10^{-4}$ \\
TRIF & $\lambda=0.01$ & $\lambda=0.02$ \\
RIF & $\lambda=0.005$ & $\lambda=0.02$ \\
RL & $K=80$ & $K=20$ \\
RLTV & $K=40,\quad \lambda=0.001$ & $K=40,\quad \lambda=0.003$ \\
\bottomrule

\end{tabular}
\end{table}

ImageJ receives the same nominal parameter values as Python. This supports a controlled transfer test. The named implementations can differ, so the scores are not an estimate of the best achievable score for a separately tuned plugin. A second comparison harmonizes the NIF floor, the RIF frequency response, and the RL update where they differ. It measures raw pixel agreement instead of maximizing a quality score.

\subsection{Quality metrics and data-fit diagnostic}

For \(Q\) aligned pixels and the fixed intensity peak \(L=1\),
\begin{equation}
\mathrm{MSE} = Q^{-1}\sum_i(\widehat f_i-f_i)^2,\qquad
\mathrm{PSNR} = 10\log_{10}(L^2/\mathrm{MSE}) .
\end{equation}
PSNR is infinite for identity. Arithmetic averages of individual PSNRs differ from PSNR computed from pooled MSE. We report per-image means, sample standard deviations, and paired differences relative to each observation. Each image has equal weight. For the auxiliary crop, MSE and PSNR use the full cropped region.

The structural similarity index compares local luminance, contrast, and structure \citep{ssim}. In the specified window,
\begin{equation}
s(f,\widehat f)=
\frac{(2\mu_f\mu_{\widehat f}+0.01^2)(2\sigma_{f\widehat f}+0.03^2)}
{(\mu_f^2+\mu_{\widehat f}^2+0.01^2)(\sigma_f^2+\sigma_{\widehat f}^2+0.03^2)} .
\end{equation}
The window is an \(11\times11\) normalized Gaussian with standard deviation 1.5, and the weighted covariance uses the population convention. The score is the mean over centers whose entire window lies in the scored region. Python computes moments through separable filtering, while Java sums the 2D weights explicitly. This gives an independent check of both the data handoff and the formula.

Raw scores use the serialized or in-memory floating-point estimate without clipping, with the same peak 1. Clipped scores use \(\clip(\widehat f)\) and \(\clip(g)\) for the observation baseline. Displays share \([0,1]\) unless a caption says otherwise. The plugin's native raw PSNR uses the reference maximum instead of 1. If that maximum is positive, native PSNR equals the fixed-peak value plus \(20\log_{10}(\max f)\). We retain both.

To distinguish statistical fit from recovery, the Poisson validation sweep also measures the per-pixel generalized KL discrepancy between the observation \(g\) and the prediction \(\widehat b=\Hoperator\widehat f\):
\begin{equation}
\frac{\KL(g\Vert\widehat b)}{Q} =
\frac1Q\sum_i\left[\widehat b_i-g_i+
 g_i\log\left(\frac{g_i}{\widehat b_i}\right)\right].
\label{eq:kl}
\end{equation}
For \(g_i=0\) the logarithmic term contributes zero. Numerical predictions at or below the positive representable floor are floored for this diagnostic. Up to constants independent of \(f\), this quantity is the normalized negative log-likelihood given in the RL derivation. It uses the raw iterate rather than its displayed version. The observed data and sharp reference remain identical at all iteration limits.

\subsection{PSF mismatch, boundary mismatch, and a geometric phantom}

The sensitivity experiment retains the Gaussian-noise settings from Table~\ref{tab:settings}. It changes the assumed Gaussian width to 1.4 or 2.4 pixels while keeping the 1.9-pixel observation. It also makes a new 9-pixel horizontal motion observation and solves it with either the correct motion PSF or the original Gaussian assumption. Finally, it generates a zero-extended blur and restores it with the circular model. The same noise realization per image is reused across the six scenarios, including the matched Gaussian baseline. Each scenario is scored on the full image and after removing 16 pixels per edge. The inner score helps localize errors without redefining the primary whole-image metric.

The geometric phantom has a known \(256\times256\) intensity array. It includes a flat background, a rectangle, a disk, point sources, two separated bars, and line groups with widths of 2, 4, and 8 pixels. It receives the same Gaussian noise scale with seed 7351 and is blurred separately with Gaussian or motion PSFs. The group did not tune a second set of settings on the phantom. Besides full-reference scores, a fixed horizontal cut through row 180 shows how the reconstructions recover or distort bar contrast.

\subsection{Independent EPFL Hollow Bars transfer experiment}

The public Hollow Bars reference and widefield PSF are \(128\times256\times256\) 16-bit stacks with ordered axes \(z,y,x\) \citep{epflbars}. We read the clean volume and PSF archives, verify their published-source identities by SHA-256, and build a two-dimensional test from the axial sums. Let \(S_z v=\sum_z v[z,\cdot,\cdot]\). For full circular convolution in all dimensions,
\begin{equation}
S_z(v *_3 k)=(S_z v)*_2(S_z k).
\label{eq:projection}
\end{equation}
To see this, exchange the sum over output \(z\) with kernel offsets. For any axial offset \(c\), the map \(z\mapsto (z-c)\bmod Z\) permutes all axial indices. A maximum-intensity projection does not share this linear identity. We compare a separate 3D FFT implementation followed by summation against the 2D convolution and obtain a maximum difference of \EpflProjectionError{}.

The axial-sum reference is scaled by its own peak, 262140. The axial-sum PSF is normalized by its total mass, 32381499, and kept on the full \(256\times256\) grid. The central \(17\times17\) block contains \EpflMassPercent\% of that mass. In contrast, the entire Gaussian PSF fits in that block. For fair intensity calibration we regenerate the two noise conditions using the same scales as BSDS500 and seeds 735101 and 735102. We use the supplied clean reference and PSF rather than the separately packaged corrupted stacks.

The projected kernel is normalized in double precision and stored as float32. In this experiment the ImageJ batch retains that PSF mass instead of renormalizing it again in Java. An initial compatibility check found that summing the wide kernel in single precision changed its scale enough to exceed the parity tolerances for the direct filters. With \file{-norm no}, both platforms use the stored normalized samples, apart from float32 roundoff. The main BSDS500 comparison continues to use plugin normalization. All reconstruction parameters are transferred unchanged from BSDS500 validation.

\section{Results}
\label{sec:results}

\subsection{Calibration on held-out images}

Figure~\ref{fig:validation} shows the mean clipped PSNR used to select parameters. For Gaussian noise, raising the RL limit from 40 to 80 iterations adds only about \(0.06\) dB. On Poisson data, the tested mean is highest at \(K=20\) and falls from \(24.35\) to \(23.63\) dB by \(K=80\). The appropriate stopping point depends on the noise and initialization. For Poisson RLTV at \(K=40\), the best mean in the weight grid occurs at \(\lambda=0.003\). Both weaker and stronger weights reduce the validation mean.

\begin{figure}[!htbp]
\centering\includegraphics[width=\linewidth]{validation.png}
\caption{Validation mean PSNR versus the method's parameter. The RLTV panels show the weight sweep at the selected iteration count. Appendix~\ref{app:validation} provides every configuration.}
\label{fig:validation}
\end{figure}

Figure~\ref{fig:rlfit} compares two objectives on the same Poisson observations. The mean per-pixel generalized KL discrepancy decreases from approximately \(7.46\times10^{-4}\) at \(K=5\) to \(6.29\times10^{-4}\) at \(K=80\), while PSNR peaks earlier. The results show that closer agreement with the noisy observation can coexist with worse recovery of the clean image. The decreasing values are empirical at the sampled iterations. They do not constitute a convergence result for the stabilized implementation.

\begin{figure}[!htbp]
\centering\includegraphics[width=\linewidth]{rl-fit.png}
\caption{Poisson validation for plain RL. Left shows clipped PSNR against ground truth. Right shows the generalized KL discrepancy between the observation and the raw iterate blurred by the known PSF. Points average the same five images.}
\label{fig:rlfit}
\end{figure}

\subsection{Natural-image reconstruction quality}

\begin{table}[!htbp]
\centering
\caption{Mean clipped quality on ten BSDS500 test images for each noise model. Python PSNR includes sample standard deviation across images.}
\label{tab:primary}
\small\setlength{\tabcolsep}{4.6pt}
\begin{tabular}{llrrrr}
\toprule
& & \multicolumn{2}{c}{Python} & \multicolumn{2}{c}{DeconvolutionLab2}\\
\cmidrule(lr){3-4}\cmidrule(lr){5-6}
Noise & Method & PSNR (dB) & SSIM & PSNR (dB) & SSIM\\\midrule
Gaussian & Observed & 23.68 $\pm$ 3.99 & 0.5990 & 23.68 & 0.5990 \\
 & NIF & 12.97 $\pm$ 0.28 & 0.1334 & 12.97 & 0.1334 \\
 & TRIF & 25.24 $\pm$ 3.91 & 0.6825 & 25.24 & 0.6825 \\
 & RIF & 25.30 $\pm$ 3.91 & 0.6825 & 25.32 & 0.6871 \\
 & RL & 25.31 $\pm$ 3.89 & 0.6844 & 25.02 & 0.6512 \\
 & RLTV & 25.44 $\pm$ 4.38 & 0.7096 & 25.14 & 0.6710 \\
Poisson & Observed & 22.97 $\pm$ 3.38 & 0.5064 & 22.97 & 0.5064 \\
 & NIF & 8.92 $\pm$ 0.50 & 0.0513 & 8.92 & 0.0513 \\
 & TRIF & 24.05 $\pm$ 3.23 & 0.6037 & 24.05 & 0.6037 \\
 & RIF & 24.36 $\pm$ 3.52 & 0.6230 & 24.40 & 0.6279 \\
 & RL & 24.51 $\pm$ 3.78 & 0.6413 & 23.56 & 0.5456 \\
 & RLTV & 24.89 $\pm$ 4.03 & 0.6684 & 23.40 & 0.5389 \\
\bottomrule

\end{tabular}
\end{table}

In Table~\ref{tab:primary}, Python RLTV yields the highest mean PSNR and SSIM among the tested settings for both noise types. Its mean gains are \(1.76\) dB for Gaussian noise and \(1.92\) dB for Poisson noise. The differences between RLTV and plain RL are approximately \(0.13\) and \(0.37\) dB respectively. These differences are smaller than the spread of absolute scores across images. Table~\ref{tab:gain} therefore pairs each result with its exact measured input. On this fixed subset, TRIF, RIF, RL, and RLTV all improve PSNR for every image and both noise types. For Python RLTV on Poisson observations, the paired improvements range from \(0.96\) to \(3.25\) dB.

NIF degrades the baseline in all main cases. Its raw Poisson mean is \(5.41\) dB, improving to only \(8.92\) dB after clipping. This gap quantifies how clipping can hide excursions in weakly transmitted bands. Since \([0,1]\) contains the reference range, componentwise clipping cannot increase squared error. It can still change local structural statistics. The paired visualization and atlas keep the same intensity scale and report the scoring policy explicitly.

\begin{table}[!htbp]
\centering
\caption{Python raw mean PSNR and paired clipped gains against the observation. SD is the sample standard deviation of the paired gain.}
\label{tab:gain}
\small\setlength{\tabcolsep}{4.2pt}
\begin{tabular}{llrrrr}
\toprule Noise & Method & Raw (dB) & Gain (dB) & SD (dB) & Improved\\\midrule
Gaussian & NIF & 12.02 & -10.72 & 3.96 & 0/10 \\
 & TRIF & 25.23 & 1.55 & 0.53 & 10/10 \\
 & RIF & 25.29 & 1.62 & 0.56 & 10/10 \\
 & RL & 25.29 & 1.63 & 0.55 & 10/10 \\
 & RLTV & 25.43 & 1.76 & 0.63 & 10/10 \\
Poisson & NIF & 5.41 & -14.05 & 3.57 & 0/10 \\
 & TRIF & 24.04 & 1.08 & 0.43 & 10/10 \\
 & RIF & 24.34 & 1.39 & 0.38 & 10/10 \\
 & RL & 24.49 & 1.54 & 0.50 & 10/10 \\
 & RLTV & 24.88 & 1.92 & 0.76 & 10/10 \\
\bottomrule

\end{tabular}
\end{table}

\begin{figure}[!htbp]
\centering\includegraphics[width=\linewidth]{paired.png}
\caption{Paired Python PSNR gain per test image and method. Each small dot represents an image and the diamond is the mean. The vertical line marks zero improvement.}
\label{fig:paired}
\end{figure}

\subsection{Verification with DeconvolutionLab2}

On the same plugin output TIFFs, maximum discrepancies between independently computed MSE, PSNR and SSIM are \MetricMSEGap, \MetricPSNRGap{} dB, and \MetricSSIMGap{}. The correction of native PSNR by the image-dependent peak agrees with fixed-peak raw PSNR to \NativePSNRGap{} dB. This confirms the measurement procedure within numerical tolerance. It leaves differences between restoration outputs to be explained by the solvers.

\begin{table}[!htbp]
\centering
\caption{Maximum absolute raw pixel difference from the plugin across all 20 main cases. Matched means that the same numerical convention was used.}
\label{tab:parity}
\small
\begin{tabular}{lrrr}
\toprule Method & Main Python & Matched Python & Tolerance\\\midrule
NIF & $2.26\times 10^{-5}$ & $2.26\times 10^{-5}$ & $1.00\times 10^{-4}$ \\
TRIF & $7.75\times 10^{-7}$ & $7.75\times 10^{-7}$ & $1.00\times 10^{-5}$ \\
RIF & $1.55\times 10^{-2}$ & $8.12\times 10^{-7}$ & $1.00\times 10^{-5}$ \\
RL & $1.59\times 10^{-1}$ & $6.14\times 10^{-6}$ & $5.00\times 10^{-5}$ \\
RLTV & $2.24\times 10^{-1}$ & --- & --- \\
\bottomrule

\end{tabular}
\end{table}

For NIF, \file{-epsilon 0.0001} sets the plugin floor on \(|H|^2\). For TRIF, both platforms use the selected \(\lambda\) and the OTF is real for the symmetric, centered PSF. For RIF, the main Python implementation uses finite differences while the plugin uses squared angular frequencies. A separate Python matching calculation uses the plugin spectrum and stabilized division. The maximum difference falls from about \(1.55\times10^{-2}\) to below \(10^{-6}\). The matching formula is verified on the even-size symmetric inputs of this study.

For RL, matching both the starting image and the division rule reduces the maximum difference from about \(1.59\times10^{-1}\) to \(6.14\times10^{-6}\). On the Poisson split at \(K=20\), the nominal Python and plugin means are \(24.51\) and \(23.56\) dB respectively. Comparing the raw arrays at matched settings establishes that the observed gap is driven by the rules being changed. RLTV also changes the discrete gradient, boundary treatment, and TV normalization. We compare its quality at the nominal settings and do not treat the outputs as pixel-equivalent.

At transferred settings the plugin RIF achieves the largest mean PSNR on the main cases, with \(25.32\) dB (Gaussian) and \(24.40\) dB (Poisson). The comparison answers how the plugin behaves on the shared experiment. The validation did not optimize its settings independently. Figure~\ref{fig:roi} and Table~\ref{tab:case} illustrate a fixed case, with the complete per-image record in Appendix~\ref{app:atlas}.

\begin{figure}[!htbp]
\centering\includegraphics[width=\linewidth]{platform-roi.png}
\caption{The identical central \(64\times64\) ROI for BSDS500 image 100007 with Poisson noise. Both rows share the fixed display range.}
\label{fig:roi}
\end{figure}

\begin{table}[!htbp]
\centering
\caption{Illustrative full-image scores for image 100007 under Poisson noise. The image panels show a smaller ROI.}
\label{tab:case}
\small\setlength{\tabcolsep}{4pt}
\begin{tabular}{lrrrrrr}
\toprule
& \multicolumn{3}{c}{Python} & \multicolumn{3}{c}{DeconvolutionLab2}\\
\cmidrule(lr){2-4}\cmidrule(lr){5-7}
& MSE \(\times10^3\) & PSNR & SSIM & MSE \(\times10^3\) & PSNR & SSIM\\\midrule
Observed & 2.327 & 26.33 & 0.5153 & 2.327 & 26.33 & 0.5153 \\
TRIF & 2.081 & 26.82 & 0.5729 & 2.081 & 26.82 & 0.5729 \\
RLTV & 1.154 & 29.38 & 0.7584 & 2.958 & 25.29 & 0.4519 \\
\bottomrule

\end{tabular}
\end{table}

\subsection{Sensitivity to PSF shape and boundary mismatch}

\begin{table}[!htbp]
\centering
\caption{Mean paired Python RLTV gain (dB) on ten Gaussian-noise test photographs. The same region is used for the observation and the estimate in each column.}
\label{tab:sens}
\small
\begin{tabular}{lrr}
\toprule Scenario & Full image & Inner region\\\midrule
Matched Gaussian & 1.76 & 1.75 \\
Assume \(\sigma\) = 1.4 & 1.22 & 1.21 \\
Assume \(\sigma\) = 2.4 & 0.86 & 0.85 \\
Matched motion & 2.96 & 2.96 \\
Motion as Gaussian & -2.79 & -2.63 \\
Zero as periodic boundary & -0.74 & 1.75 \\
\bottomrule

\end{tabular}
\end{table}

As shown in Table~\ref{tab:sens}, moving the assumed Gaussian width from 1.9 to 1.4 or 2.4 pixels reduces mean RLTV gain from \(1.76\) to \(1.22\) or \(0.86\) dB. On motion observations, using the correct 9-pixel box gives a \(2.96\) dB gain, whereas assuming the Gaussian PSF incurs a \(2.79\) dB loss. Removing the border leaves a \(2.63\) dB loss for the wrong blur model. That loss persists away from the edge because the forward model is wrong throughout the image.

The zero-boundary case behaves differently. Restoring it with the periodic model loses \(0.74\) dB over the full frame but gains \(1.75\) dB in the cropped interior. This does not justify excluding boundaries from the primary comparison. It demonstrates where the model error is concentrated. Figure~\ref{fig:sensitivity} shows corresponding trends for the other regularized and iterative solvers. The full CSV also retains NIF, whose much lower gains would dominate the plot scale.

\begin{figure}[!htbp]
\centering\includegraphics[width=\linewidth]{sensitivity.png}
\caption{Mean paired PSNR change for regularized and iterative methods under the six sensitivity scenarios. Both panels use the frozen Gaussian-noise configuration.}
\label{fig:sensitivity}
\end{figure}

\subsection{Geometric phantom and line contrast}

In the Gaussian phantom experiment, Python RIF and RLTV obtain \(21.66\) and \(21.63\) dB. RIF has SSIM \(0.7285\), below the observation's \(0.7611\), while RLTV obtains \(0.8384\). The similar PSNRs hide a structural difference on a target with extensive flat regions. In the motion case, all methods except NIF improve PSNR, yet every method falls below the observation in SSIM. Plain RL, for example, reaches \(23.12\) dB and \(0.4414\) SSIM. The result is one geometric configuration, with parameters transferred from the photographs.

\begin{figure}[!htbp]
\centering\includegraphics[width=\linewidth]{phantom.png}
\caption{Geometric reference, blurred observations and selected Python restorations. Top uses Gaussian blur and bottom uses horizontal motion blur. The phantom is identical in both rows.}
\label{fig:phantom}
\end{figure}

The profile in Figure~\ref{fig:profile} distinguishes the isolated bright pair from groups of narrow repetitive lines. A broad kernel smooths their contrast differently. Periodic motion blur has notches in its transfer function, so tuning the inverse cannot restore all spatial frequencies equally. The lines, their background, and associated oscillations remain visible when the profile is scored at its raw intensities. Appendix~\ref{app:additional} reports all phantom scores.

\begin{figure}[!htbp]
\centering\includegraphics[width=.96\linewidth]{profile.png}
\caption{Fixed raw-intensity cross-section through the phantom at row 180. The displayed range of the image panels does not clip these curves.}
\label{fig:profile}
\end{figure}

\subsection{Transfer to the EPFL optical PSF}

\begin{table}[!htbp]
\centering
\caption{Clipped quality on the two EPFL Hollow Bars projection cases. There is one observation per noise model, with settings transferred from BSDS500.}
\label{tab:epfl}
\small\setlength{\tabcolsep}{4.5pt}
\begin{tabular}{llrrrr}
\toprule
& & \multicolumn{2}{c}{Python} & \multicolumn{2}{c}{DeconvolutionLab2}\\
\cmidrule(lr){3-4}\cmidrule(lr){5-6}
Noise & Method & PSNR & SSIM & PSNR & SSIM\\\midrule
Gaussian & Observed & 13.19 & 0.1162 & 13.19 & 0.1162 \\
 & NIF & 8.38 & 0.0785 & 8.38 & 0.0785 \\
 & TRIF & 14.77 & 0.5133 & 14.77 & 0.5133 \\
 & RIF & 17.69 & 0.3929 & 17.61 & 0.3965 \\
 & RL & 15.10 & 0.6369 & 15.42 & 0.6479 \\
 & RLTV & 14.30 & 0.5747 & 14.73 & 0.5913 \\
Poisson & Observed & 13.19 & 0.1367 & 13.19 & 0.1367 \\
 & NIF & 9.42 & 0.0485 & 9.42 & 0.0485 \\
 & TRIF & 14.15 & 0.5659 & 14.15 & 0.5659 \\
 & RIF & 16.70 & 0.5627 & 16.66 & 0.5636 \\
 & RL & 13.82 & 0.5207 & 14.23 & 0.5444 \\
 & RLTV & 14.24 & 0.5867 & 14.70 & 0.6054 \\
\bottomrule

\end{tabular}
\end{table}

The EPFL projection substantially changes the distribution of transmitted frequencies and the pixel intensity histogram. Python RIF achieves \(17.69\) and \(16.70\) dB, improving over the observations by about \(4.50\) and \(3.51\) dB. It has the highest PSNR under both noise types. The highest Python SSIM belongs to RL in the Gaussian case and RLTV in the Poisson case. The different rankings show that the BSDS500 result for RLTV does not transfer uniformly to the optical PSF and sparse bars.

Plugin RIF also has the highest PSNR for both EPFL cases. With mass preserved and the matched rules applied, the maximum raw difference over NIF, TRIF, RIF, and RL is \EpflParityError{}. The arithmetic agreement on a broad kernel confirms the centering and normalization strategy beyond the Gaussian kernel. It does not validate a three-dimensional restoration. Figure~\ref{fig:epfl} visualizes the identical 2D target and selected outputs. The broad PSF and shared display window make the unprocessed observations faint.

\begin{figure}[!htbp]
\centering\includegraphics[width=\linewidth]{epfl.png}
\caption{Axial-sum EPFL Hollow Bars projection and selected Python restorations. Top row uses Gaussian noise and bottom row uses Poisson noise. All images use the same \([0,1]\) display range.}
\label{fig:epfl}
\end{figure}

\subsection{Computational performance}

\begin{table}[!htbp]
\centering
\caption{Python solver timing on validation image 101085 with Poisson noise. Values are medians of seven alternating measurements after warm-up.}
\label{tab:performance}
\small\setlength{\tabcolsep}{4pt}
\begin{tabular}{lrrrrrr}
\toprule & \multicolumn{2}{c}{Time (ms)} && \multicolumn{2}{c}{Tracked MiB} &\\
Method & Full FFT & Real FFT & Ratio & Full FFT & Real FFT & Max diff.\\\midrule
NIF & 4.28 & 2.42 & 1.77 & 6.63 & 3.65 & $1.78\times 10^{-15}$ \\
TRIF & 4.91 & 2.39 & 2.05 & 6.13 & 3.40 & $3.89\times 10^{-16}$ \\
RIF & 4.43 & 2.30 & 1.92 & 7.13 & 3.91 & $4.44\times 10^{-16}$ \\
RL & 116.17 & 38.66 & 3.01 & 7.00 & 4.52 & $2.89\times 10^{-15}$ \\
RLTV & 293.73 & 156.45 & 1.88 & 10.00 & 8.51 & $2.11\times 10^{-15}$ \\
\bottomrule

\end{tabular}
\end{table}

Run \texttt{\RunID} used Python \PythonVersion, NumPy \NumPyVersion, and SciPy \SciPyVersion{} on Windows. Real FFT and OTF caching yield a measured speed ratio of \BenchmarkSpeedRange{}. The largest numeric difference between the two branches is \BenchmarkMaxError{}. Array sizes and work per iteration explain why iterative methods spend more time than closed-form filters. TV also computes spatial differences and normalization, which are not accelerated by FFT reuse. Since the benchmark changes storage format and caching together, the table does not isolate their individual effects. Tracked allocation is not process RSS. The benchmark is specific to one input, solver configuration, and environment.

\FloatBarrier
\section{Discussion}
\label{sec:discussion}

\subsection{Interpreting improvement}

The two Tikhonov variants reduce the sensitivity that makes NIF unstable. RL gains quality at finite iteration limits, and adding TV provides additional benefit on the main photographed scenes at the selected settings. The Poisson KL diagnostic separates data fit from recovery: after the oracle PSNR peak, the objective in the measurement domain continues to decrease. In a practical deployment the ground truth is rarely available, so an iteration rule tuned on representative validation data is more realistic than maximizing test PSNR.

Different methods can lead on different targets or metrics. The geometric phantom and EPFL projection demonstrate this directly. A regularizer that helps a smooth background can alter edges or narrow lines. Conversely, a method can sharpen the bars enough to reduce squared error while creating weak structured artifacts. The study therefore uses absolute and paired PSNR, local SSIM, intensity profiles, and common-scale images. Additional resolution claims, such as resolving two physical emitters, would need a calibrated spatial experiment and a resolution criterion.

\subsection{Distinguishing algorithm and implementation}

Matching input TIFFs is necessary but insufficient for comparing software. A normalization error on a broad PSF, a different definition of the Laplacian, or one pixel of phase shift changes the reconstructed array. Tolerances are applied only after defining the compared discrete operators. The border impulse, raw pixel errors, and independent score implementations give three complementary checks. ImageJ is a reference implementation, and it need not use the same default choices as the educational code.

RIF and RL illustrate the distinction. The main RIF uses the five-point operator. The matching branch uses the plugin's spectral convention. Plain RL updates with two different initializations and numerical ratio rules because the test evaluates both the project configuration and like-for-like agreement. RLTV uses even more implementation-specific TV details. Presenting all those values under one unlabeled column would conflate numeric verification with a comparison of method choices.

\subsection{Practical failure diagnosis}

Ringing near a contour suggests a frequency gain that is too large or an incorrect boundary model. A global image shift suggests an origin mismatch. Noise amplified over a flat region suggests weak regularization or too many iterations. A strong color LUT or ImageJ auto-contrast can make different float arrays appear similar. In the present workflow the primary diagnostics are the centered PSF, spatial impulse, raw TIFF, and numerical errors. A visual comparison always shares a display range. The paired boundary experiment shows why inspecting the margins is important even when the center improves.

A practical choice among these methods depends on the expected noise, the kernel, and the task. Identity-penalty Tikhonov gives a compact direct baseline. Laplacian regularization preserves DC and penalizes local oscillation. Plain RL is appropriate when photon statistics and nonnegativity matter and a stopping rule is available. RLTV can suppress noise on piecewise-smooth structures, with increased cost and a stronger dependence on the TV convention and weight. No line of Table~\ref{tab:primary} establishes a universal ranking.

\subsection{Limitations and external validity}

The natural-image test is a fixed 10-image subset and uses one noise realization for each image and noise model. Cross-image standard deviations are descriptive. They do not capture variation across repeated draws, calibrated cameras, or the full BSDS500 distribution. The simulated Gaussian noise is clipped below zero and the Poisson generator operates on normalized, gamma-encoded JPEG crops. Neither measurement pipeline fully models a physical sensor.

The principal experiments know the blur PSF and use a shared periodic convention. The mismatch experiments explore selected deviations, including a linear zero-extended image, without estimating the correct model from data. The EPFL experiment is one independently sourced synthetic structure projected to 2D. It does not establish performance on acquired microscopy, mixed detector noise, unknown backgrounds, spatially varying blur, color channels, or full 3D reconstruction. More expansive claims would require additional datasets with independent calibration, multiple noise draws, and separate validation for each platform.

\section{Conclusion}

We developed the non-blind 2D forward model, derived five direct and iterative estimators, implemented them in Python, and verified the shared computation against DeconvolutionLab2 in ImageJ. Stabilized inverse filtering remained sensitive to noise. Tikhonov and finite-step Richardson--Lucy variants improved the principal observations, and RLTV led the Python mean scores at our frozen BSDS500 settings. The EPFL and phantom experiments showed that the best method and metric can change with structure and PSF. The numerical evidence includes identical test inputs, independent scores, raw output comparisons, and an origin-sensitive impulse check. The complete case atlas and scripts permit independent examination of these conclusions.

\FloatBarrier
\begin{thebibliography}{20}
\raggedright
\bibitem[Arbel\'aez et al.(2011)Arbel\'aez, Maire, Fowlkes, and Malik]{bsds}
P. Arbel\'aez, M. Maire, C. Fowlkes, and J. Malik.
Contour detection and hierarchical image segmentation.
\emph{IEEE Transactions on Pattern Analysis and Machine Intelligence}, 33(5):898--916, 2011.
\href{https://doi.org/10.1109/TPAMI.2010.161}{doi:10.1109/TPAMI.2010.161}.

\bibitem[Berkeley Computer Vision Group(2013)]{bsdsdata}
Berkeley Computer Vision Group.
\emph{Berkeley Segmentation Data Set and Benchmarks 500}.
\url{https://www2.eecs.berkeley.edu/Research/Projects/CS/vision/grouping/resources.html}.

\bibitem[Dey et al.(2006)Dey, Blanc-F\'eraud, Zimmer, Roux, Kam, Olivo-Marin, and Zerubia]{dey}
N. Dey, L. Blanc-F\'eraud, C. Zimmer, P. Roux, Z. Kam, J.-C. Olivo-Marin, and J. Zerubia.
Richardson--Lucy algorithm with total variation regularization for 3D confocal microscope deconvolution.
\emph{Microscopy Research and Technique}, 69(4):260--266, 2006.
\href{https://doi.org/10.1002/jemt.20294}{doi:10.1002/jemt.20294}.

\bibitem[EPFL Biomedical Imaging Group(2018)]{dl2docs}
EPFL Biomedical Imaging Group.
\emph{DeconvolutionLab2 documentation and version 2.1.2 distribution}, 2018.
\url{https://bigwww.epfl.ch/deconvolution/deconvolutionlab2/}.

\bibitem[EPFL Biomedical Imaging Group(n.d.-a)]{dl2code}
EPFL Biomedical Imaging Group.
\emph{DeconvolutionLab2 source code}.
Source comparison at commit \file{e9af0aba493ba137d70877154648e5583e376a81}.
\url{https://github.com/Biomedical-Imaging-Group/DeconvolutionLab2}.

\bibitem[EPFL Biomedical Imaging Group(n.d.-b)]{epflbars}
EPFL Biomedical Imaging Group.
\emph{Synthetic hollow bars}, public clean stack and widefield PSF.
\url{https://bigwww.epfl.ch/deconvolution/bars/}.

\bibitem[Fessler(2017)]{fessler}
J. A. Fessler.
Regularization.
Chapter 2 of \emph{Image Reconstruction}, University of Michigan, 2017.
\url{https://web.eecs.umich.edu/~fessler/book/c-reg.pdf}.

\bibitem[Hansen et al.(2006)Hansen, Nagy, and O'Leary]{hansen}
P. C. Hansen, J. G. Nagy, and D. P. O'Leary.
\emph{Deblurring Images: Matrices, Spectra, and Filtering}.
SIAM, 2006.
\href{https://doi.org/10.1137/1.9780898718874}{doi:10.1137/1.9780898718874}.

\bibitem[Kirshner and Sage(2017)]{psfgen}
H. Kirshner and D. Sage.
\emph{PSF Generator}, EPFL software and model documentation, 2017 distribution.
\url{https://bigwww.epfl.ch/algorithms/psfgenerator/}.

\bibitem[Lucy(1974)]{lucy}
L. B. Lucy.
An iterative technique for the rectification of observed distributions.
\emph{The Astronomical Journal}, 79:745--754, 1974.
\href{https://doi.org/10.1086/111605}{doi:10.1086/111605}.

\bibitem[Richardson(1972)]{richardson}
W. H. Richardson.
Bayesian-based iterative method of image restoration.
\emph{Journal of the Optical Society of America}, 62(1):55--59, 1972.
\href{https://doi.org/10.1364/JOSA.62.000055}{doi:10.1364/JOSA.62.000055}.

\bibitem[Sage et al.(2017)Sage, Donati, Soulez, Fortun, Schmit, Seitz, Guiet, Vonesch, and Unser]{sage}
D. Sage, L. Donati, F. Soulez, D. Fortun, G. Schmit, A. Seitz, R. Guiet, C. Vonesch, and M. Unser.
DeconvolutionLab2: An open-source software for deconvolution microscopy.
\emph{Methods}, 115:28--41, 2017.
\href{https://doi.org/10.1016/j.ymeth.2016.12.015}{doi:10.1016/j.ymeth.2016.12.015}.

\bibitem[Shepp and Vardi(1982)]{shepp}
L. A. Shepp and Y. Vardi.
Maximum likelihood reconstruction for emission tomography.
\emph{IEEE Transactions on Medical Imaging}, 1(2):113--122, 1982.
\href{https://doi.org/10.1109/TMI.1982.4307558}{doi:10.1109/TMI.1982.4307558}.

\bibitem[Tikhonov(1963)]{tikhonov}
A. N. Tikhonov.
On the solution of ill-posed problems and the method of regularization.
\emph{Doklady Akademii Nauk SSSR}, 151(3):501--504, 1963.
\url{https://www.mathnet.ru/eng/dan/v151/i3/p501}.

\bibitem[Wang et al.(2004)Wang, Bovik, Sheikh, and Simoncelli]{ssim}
Z. Wang, A. C. Bovik, H. R. Sheikh, and E. P. Simoncelli.
Image quality assessment: From error visibility to structural similarity.
\emph{IEEE Transactions on Image Processing}, 13(4):600--612, 2004.
\href{https://doi.org/10.1109/TIP.2003.819861}{doi:10.1109/TIP.2003.819861}.
\end{thebibliography}

\clearpage
\appendix
\section{Reproduction and artifact inventory}
\label{app:reproduce}

\subsection{Environment and data}

The report uses run \texttt{\RunID}, produced with Python \PythonVersion, NumPy \NumPyVersion, SciPy \SciPyVersion, ImageJ 1.54p, and Java \JVMVersion. Python dependencies are pinned in \file{code/requirements.txt}. The manifests record SHA-256 hashes for the Berkeley and EPFL data and the ImageJ and DeconvolutionLab2 JARs.

\textbf{Validation JPEG IDs}\quad\ValidationIDs.

\textbf{Test JPEG IDs}\quad\TestIDs.

\subsection{Commands}

Run Listing~\ref{lst:commands} from the project root with \file{java} and \file{javac} on the path. The follow-up experiments read the latest BSDS500 run and reuse its frozen parameters.

\begin{lstlisting}[style=paper,language=bash,caption={Commands for the experiments and numerical checks.},
label={lst:commands}]
python code/scripts/prepare_data.py
python code/scripts/check.py
python code/scripts/run_experiments.py
python code/scripts/audit_artifacts.py
python code/scripts/run_imagej.py
python code/scripts/run_sensitivity.py
python code/scripts/run_epfl_bars.py
\end{lstlisting}

\clearpage
\subsection{Expected artifacts and audit counts}

\begin{center}
\captionof{table}{Checklist of numerical evidence. The paths are relative to a completed run directory.}
\small
\begin{tabularx}{\linewidth}{@{}p{.44\linewidth}Y@{}}
\toprule Artifact & Expected content\\\midrule
\file{config_frozen_before_test.json} & Selected values and candidate grids recorded before any test reconstruction.\\
\file{validation_sweep.csv} & One row per method setting, validation image, and noise model.\\
\file{validation_rl_fit.csv} & 25 records of RL PSNR and observation-domain KL.\\
\file{test_metrics.csv} & 240 rows for 10 photographs, two noise models, six methods including the observation, and two scoring policies.\\
\file{artifact_audit.json} & Python TIFF coverage and metric recomputation.\\
\file{imagej/imagej_metrics.csv} & 242 rows including both scores for the asymmetric impulse output.\\
\file{imagej/parity.csv} & Raw pixel comparisons for 100 restorations.\\
\file{sensitivity/metrics.csv} & 720 rows across six scenarios, ten photos, six methods, and two regions.\\
\file{sensitivity/phantom_metrics.csv} & 12 rows from two known PSFs and six method outputs each.\\
\file{epfl-bars/manifest.json} & Archive hashes, geometry, seeds, mass fraction, and 3D-to-2D error.\\
\file{epfl-bars/imagej/audit.json} & 26 plugin metric rows, including the two impulse scores.\\
\file{benchmark.csv} & Five timing ratios, raw repetitions, and maximum numeric differences.\\
\bottomrule
\end{tabularx}
\end{center}

Each image and noise case retains its reference, observation, noise-free blur, PSF, and raw estimates as float TIFF. The audits recompute scores from the saved pixels.

\subsection{Team responsibilities}

\begin{tabularx}{\linewidth}{@{}p{.29\linewidth}Y@{}}
\toprule Student & Technical focus\\\midrule
Nguyen Minh Hieu & Forward model, boundary assumptions, Fourier derivation, and worked examples.\\
Pham Manh Dung & Python solvers, optimization, and unit testing.\\
Nguyen Thanh Tam & DeconvolutionLab2 inspection, ImageJ workflow, macros, and Python CLI handoff for reference runs.\\
Vu Dinh Hieu & Data preparation, evaluation protocol, metric aggregation, and sensitivity experiments.\\
\bottomrule
\end{tabularx}

\clearpage
\section{Independent numerical oracles}
\label{app:oracles}

\subsection{Impulse response and asymmetric adjoint}

Convolution with a unit impulse places a copy of the PSF at the impulse origin. The ImageJ probe uses an impulse at zero-based pixel \((0,31)\) in a \(32\times32\) array and a \(3\times3\) asymmetric kernel:
\[
h=\begin{pmatrix}
0&0.1&0.2\\
0.3&0.4&0\\
0&0&0
\end{pmatrix}.
\]
The nonzero output offsets place values \(0.1,0.2,0.3,0.4\) at pixels
\((31,31),(31,0),(0,30),(0,31)\), respectively. A correlation reverses their spatial orientation. A zero extension drops mass outside the image. A center error shifts the entire pattern. Comparing the full float array therefore tests all three failure modes at once. The independent Python solver tests also use spatial convolution without FFT when calculating a known asymmetric forward or adjoint result.

\subsection{Spectral and dense least-squares agreement}

For a small image, construct the columns of \(\Hoperator\) by applying spatial convolution to each unit basis vector in turn. Construct \(L\) the same way using the Laplacian stencil. A dense solve of
\((\Hoperator^\top\Hoperator+\lambda L^\top L)x=\Hoperator^\top g\)
then checks the spectral division and multiplication by \(\overline H\) without using the implementation's OTF. The test also uses \(L=I\) for TRIF. This comparison distinguishes a correct periodic estimator from a plausible-looking but shifted restoration. A one-row identity-blur case verifies when the vertical neighbor terms cancel.

\subsection{Hand verification of one likelihood update}

For the two-pixel circular operator with diagonal \(3/4\), take \(g=(6,2)\) and \(f^{(0)}=(4,4)\). Both prediction values are 4. The back-projected ratio is therefore \((1.25,0.75)\), which gives \(f^{(1)}=(5,3)\). The next prediction is \((4.5,3.5)\). Its ratio is \((4/3,4/7)\). The two back-projected components are \(8/7\) and \(16/21\), yielding the second iterate \((40/7,16/7)\). Both iterates sum to the observed total 8. This also validates the use of the adjoint in the update rather than a direct additive correction.

\subsection{Directional verification of total variation}

Write \(E(f)=\operatorname{TV}_\delta(f)\). For a deterministic nonconstant test image and perturbation \(w\), compare
\[
\frac{E(f+\eta w)-E(f-\eta w)}{2\eta}
\quad\text{and}\quad
\left\langle-\operatorname{div}
\left(\frac{D_+ f}{\sqrt{\norm{D_+ f}^2+\delta^2}}\right),w\right\rangle .
\]
This verifies the sign of the discrete divergence and the normalization on the same grid. Further checks compare RL and RLTV at \(\lambda=0\), preserve a constant image to numerical tolerance, and ensure the all-zero input stays zero. The positivity condition is checked in addition to the empirical tests.

\subsection{Metric and I/O checks}

An \(11\times11\) example has exactly one valid SSIM center. Its weighted mean, variance, and cross covariance are calculated explicitly from the window and compared to the filtered-moment implementation. Identity must give SSIM 1 and infinite PSNR. An error of 0.1 at every pixel gives MSE 0.01 and 20 dB with peak 1. Float TIFF round-trip tests include values outside \([0,1]\) to ensure that export does not convert a raw estimate into an 8-bit display. The application-level audit then checks every expected experiment row rather than just the number of rows.

\clearpage
\section{ImageJ runbook and comparison settings}
\label{app:imagej}

\subsection{Manual case}

The case \file{100007_poisson} is a useful replay. Open its \file{observed.tif} and the comparison folder's \file{psf_centered.tif}. The first is a 32-bit single-slice image. The second is an image-sized PSF with origin \((128,128)\); it is not the small Python PSF displayed with contrast stretching. Select TRIF and enter \(\lambda=0.02\). Set lateral and axial padding and apodization to none. Select Academic FFT, unit PSF normalization, and no intensity constraint. Export an intact float stack. Select the original case's \file{reference.tif} for scoring, preserving image alignment and float units.

\begin{lstlisting}[style=paper,caption={Option string for the main TRIF example. Paths stand for the shared float TIFFs.},
label={lst:ijargs}]
-image file OBSERVED.tif -psf file PSF_CENTERED.tif
-algorithm TRIF 0.02 -epsilon 0.000001
-pad NO NO 0 0 -apo NO NO -norm 1 -fft Academic
-monitor console -stats false -constraint no
-multithreading no -path RESULT_DIRECTORY
-out stack TRIF intact float noshow
\end{lstlisting}

The Java launcher passes the same options as one string to the plugin's entry point after setting the active macro options. \file{-stats false} and \file{noshow} avoid GUI-dependent presentation during batch operation. The console log records the full command and the start and end of each solver. An exception or incomplete output is treated as a failed run rather than a successful image.

\subsection{Auditing agreement}

\begin{tabularx}{\linewidth}{@{}lYY@{}}
\toprule Check & Shared input or setting & Interpretation\\\midrule
TIFF decoding & Exact reference and result arrays & Agreement isolates the metric routine.\\
Impulse & Border position and asymmetric kernel & Measures origin, orientation, and wrap.\\
NIF & Same floor in both FFT implementations & Agreement tests inverse gain and phase.\\
TRIF & Shared weight and symmetric center & Agreement tests the identity-penalty solution.\\
RIF & Matching auxiliary spectral Laplacian & Separates discrete operators from the solver.\\
RL & Matching initialization and ratio rule & Separates iteration conventions from FFT.\\
RLTV & Same nominal parameters only & Scores measure two distinct TV variants.\\
\bottomrule
\end{tabularx}

In the EPFL experiment the common PSF is already normalized in Python, so the option changes from \file{-norm 1} to \file{-norm no}. The run retains the full PSF support instead of forcing it to fit the much smaller Gaussian kernel. The protocol only generalizes agreement to the recorded 2D kernels and algorithm variants.

\clearpage
\section{Complete validation sweep}
\label{app:validation}

The validation performance below averages the same five images in each noise model. NIF is fixed and has a single entry. Frozen choices are marked with an asterisk, including the sole NIF configuration. RLTV varies \(K\) and \(\lambda\) jointly while holding \(\delta=0.01\). PSNR ties displayed after rounding are resolved using the underlying unrounded means.

\begin{center}
\small
\captionof{table}{Complete set of Gaussian-noise hyperparameter candidates.}
\begin{tabular}{llr}
\toprule Method & Candidate & Validation PSNR (dB)\\\midrule
NIF & $\tau=10^{-4}$ * & 13.028 \\
TRIF & $\lambda=0.0001$ & 15.447 \\
TRIF & $\lambda=0.001$ & 22.760 \\
TRIF & $\lambda=0.01$ * & 24.951 \\
TRIF & $\lambda=0.02$ & 24.741 \\
RIF & $\lambda=0.0001$ & 20.362 \\
RIF & $\lambda=0.001$ & 24.380 \\
RIF & $\lambda=0.005$ * & 25.017 \\
RIF & $\lambda=0.02$ & 24.864 \\
RL & $K=5$ & 23.971 \\
RL & $K=10$ & 24.391 \\
RL & $K=20$ & 24.726 \\
RL & $K=40$ & 24.957 \\
RL & $K=80$ * & 25.018 \\
RLTV & $K=10,\quad \lambda=0.0001$ & 24.391 \\
RLTV & $K=10,\quad \lambda=0.001$ & 24.383 \\
RLTV & $K=10,\quad \lambda=0.003$ & 24.359 \\
RLTV & $K=10,\quad \lambda=0.01$ & 24.254 \\
RLTV & $K=20,\quad \lambda=0.0001$ & 24.727 \\
RLTV & $K=20,\quad \lambda=0.001$ & 24.725 \\
RLTV & $K=20,\quad \lambda=0.003$ & 24.695 \\
RLTV & $K=20,\quad \lambda=0.01$ & 24.525 \\
RLTV & $K=40,\quad \lambda=0.0001$ & 24.965 \\
RLTV & $K=40,\quad \lambda=0.001$ * & 24.995 \\
RLTV & $K=40,\quad \lambda=0.003$ & 24.955 \\
RLTV & $K=40,\quad \lambda=0.01$ & 24.698 \\
\bottomrule

\end{tabular}
\end{center}

\clearpage
\begin{center}
\small
\captionof{table}{Complete set of Poisson-noise hyperparameter candidates.}
\begin{tabular}{llr}
\toprule Method & Candidate & Validation PSNR (dB)\\\midrule
NIF & $\tau=10^{-4}$ * & 8.919 \\
TRIF & $\lambda=0.0001$ & 10.512 \\
TRIF & $\lambda=0.001$ & 18.145 \\
TRIF & $\lambda=0.01$ & 23.672 \\
TRIF & $\lambda=0.02$ * & 24.071 \\
RIF & $\lambda=0.0001$ & 14.977 \\
RIF & $\lambda=0.001$ & 21.212 \\
RIF & $\lambda=0.005$ & 23.625 \\
RIF & $\lambda=0.02$ * & 24.300 \\
RL & $K=5$ & 23.893 \\
RL & $K=10$ & 24.216 \\
RL & $K=20$ * & 24.348 \\
RL & $K=40$ & 24.198 \\
RL & $K=80$ & 23.625 \\
RLTV & $K=10,\quad \lambda=0.0001$ & 24.217 \\
RLTV & $K=10,\quad \lambda=0.001$ & 24.222 \\
RLTV & $K=10,\quad \lambda=0.003$ & 24.224 \\
RLTV & $K=10,\quad \lambda=0.01$ & 24.164 \\
RLTV & $K=20,\quad \lambda=0.0001$ & 24.355 \\
RLTV & $K=20,\quad \lambda=0.001$ & 24.406 \\
RLTV & $K=20,\quad \lambda=0.003$ & 24.456 \\
RLTV & $K=20,\quad \lambda=0.01$ & 24.394 \\
RLTV & $K=40,\quad \lambda=0.0001$ & 24.229 \\
RLTV & $K=40,\quad \lambda=0.001$ & 24.423 \\
RLTV & $K=40,\quad \lambda=0.003$ * & 24.592 \\
RLTV & $K=40,\quad \lambda=0.01$ & 24.535 \\
\bottomrule

\end{tabular}
\end{center}

\begin{center}
\small
\captionof{table}{Plain RL data fitting and reference recovery on the five Poisson validation cases.}
\begin{tabular}{rrr}
\toprule Iterations & Mean PSNR (dB) & Mean KL per pixel\\\midrule
5 & 23.893 & $7.46\times 10^{-4}$ \\
10 & 24.216 & $6.89\times 10^{-4}$ \\
20 & 24.348 & $6.61\times 10^{-4}$ \\
40 & 24.198 & $6.43\times 10^{-4}$ \\
80 & 23.625 & $6.29\times 10^{-4}$ \\
\bottomrule

\end{tabular}
\end{center}

\clearpage
\section{Additional quantitative results}
\label{app:additional}

\subsection{Complete phantom scores}

The geometric experiments share the same reference and noise realization but use different forward PSFs. The figures in the body show selected estimators so the mark patterns remain legible. The table includes the observation and all five estimators. A strongly improved mean squared error need not produce the highest structural similarity when a reconstruction changes contrast in narrow repeated bars.

\begin{center}
\small
\captionof{table}{Geometric phantom quality with matched PSFs. Metrics use clipped intensities.}
\begin{tabular}{llrr}
\toprule PSF & Method & PSNR (dB) & SSIM\\\midrule
Gaussian & Observed & 19.28 & 0.7611 \\
Gaussian & NIF & 13.97 & 0.1714 \\
Gaussian & TRIF & 21.58 & 0.7450 \\
Gaussian & RIF & 21.66 & 0.7285 \\
Gaussian & RL & 21.66 & 0.7418 \\
Gaussian & RLTV & 21.63 & 0.8384 \\
Motion & Observed & 17.37 & 0.7120 \\
Motion & NIF & 14.09 & 0.1882 \\
Motion & TRIF & 21.76 & 0.4348 \\
Motion & RIF & 23.06 & 0.6536 \\
Motion & RL & 23.12 & 0.4414 \\
Motion & RLTV & 22.77 & 0.6290 \\
\bottomrule

\end{tabular}
\end{center}

\subsection{Sensitivity across estimators}

The following full-frame gains are relative to each scenario's own observation. A negative value indicates that recovery is less accurate by MSE than leaving the observed image as is. All rows reuse the same Gaussian-noise configuration rather than selecting favorable parameters on the test images.

\begin{center}
\small\setlength{\tabcolsep}{4pt}
\captionof{table}{Mean full-image PSNR gain (dB) for all methods and sensitivity scenarios.}
\begin{tabular}{lrrrrr}
\toprule Scenario & NIF & TRIF & RIF & RL & RLTV\\\midrule
Matched Gaussian & -10.72 & 1.55 & 1.62 & 1.63 & 1.76 \\
Assume \(\sigma\) = 1.4 & -12.81 & 0.50 & 1.02 & 0.59 & 1.22 \\
Assume \(\sigma\) = 2.4 & -9.81 & 0.96 & 0.08 & 0.47 & 0.86 \\
Matched motion & -11.36 & 0.47 & 2.74 & 0.16 & 2.96 \\
Motion as Gaussian & -15.76 & -4.15 & -4.85 & -4.71 & -2.79 \\
Zero as periodic boundary & -11.05 & -0.76 & -1.05 & -1.40 & -0.74 \\
\bottomrule

\end{tabular}
\end{center}

\clearpage
\section{Per-image test atlas}
\label{app:atlas}

Each page covers one of the ten test photographs under both noise models. The first two rows show full frames and the last two show fixed central \(64\times64\) crops. All panels use the \([0,1]\) display range. The tables score all five methods on the full frames. Gains are paired with the corresponding observation. The CSV snapshot retains MSE, PSNR and SSIM for both raw and clipped estimates.



\subsection{BSDS500 image 100007}
\begin{center}
\includegraphics[width=.92\linewidth]{atlas-100007.png}
\captionof{figure}{Full-frame and central-ROI comparison for image 100007.}
\end{center}
\begin{center}
\small\setlength{\tabcolsep}{3.3pt}
\captionof{table}{Clipped whole-image results. Gain is relative to the observation (dB).}
\begin{tabular}{llrrrrrr}
\toprule
& & \multicolumn{3}{c}{Python} & \multicolumn{3}{c}{DeconvolutionLab2}\\
\cmidrule(lr){3-5}\cmidrule(lr){6-8}
Noise & Method & PSNR & SSIM & Gain & PSNR & SSIM & Gain\\\midrule
Gaussian & Observed & 28.25 & 0.7385 & 0.00 & 28.25 & 0.7385 & 0.00 \\
 & NIF & 13.46 & 0.0754 & -14.78 & 13.46 & 0.0754 & -14.78 \\
 & TRIF & 29.67 & 0.7482 & 1.43 & 29.67 & 0.7482 & 1.43 \\
 & RIF & 29.81 & 0.7407 & 1.57 & 29.92 & 0.7524 & 1.67 \\
 & RL & 29.79 & 0.7393 & 1.54 & 29.11 & 0.6784 & 0.87 \\
 & RLTV & 30.57 & 0.8272 & 2.33 & 29.58 & 0.7307 & 1.34 \\
Poisson & Observed & 26.33 & 0.5153 & 0.00 & 26.33 & 0.5153 & 0.00 \\
 & NIF & 8.15 & 0.0214 & -18.18 & 8.15 & 0.0214 & -18.18 \\
 & TRIF & 26.82 & 0.5729 & 0.48 & 26.82 & 0.5729 & 0.48 \\
 & RIF & 27.68 & 0.6103 & 1.35 & 27.86 & 0.6250 & 1.53 \\
 & RL & 28.30 & 0.6631 & 1.97 & 25.90 & 0.4705 & -0.43 \\
 & RLTV & 29.38 & 0.7584 & 3.05 & 25.29 & 0.4519 & -1.04 \\
\bottomrule

\end{tabular}
\end{center}


\clearpage
\subsection{BSDS500 image 100039}
\begin{center}
\includegraphics[width=\linewidth]{atlas-100039.png}
\captionof{figure}{Full-frame and central-ROI comparison for image 100039.}
\end{center}
\begin{center}
\small\setlength{\tabcolsep}{3.3pt}
\captionof{table}{Clipped whole-image results. Gain is relative to the observation (dB).}
\begin{tabular}{llrrrrrr}
\toprule
& & \multicolumn{3}{c}{Python} & \multicolumn{3}{c}{DeconvolutionLab2}\\
\cmidrule(lr){3-5}\cmidrule(lr){6-8}
Noise & Method & PSNR & SSIM & Gain & PSNR & SSIM & Gain\\\midrule
Gaussian & Observed & 20.71 & 0.5292 & 0.00 & 20.71 & 0.5292 & -0.00 \\
 & NIF & 12.76 & 0.1600 & -7.95 & 12.76 & 0.1600 & -7.95 \\
 & TRIF & 22.11 & 0.6626 & 1.40 & 22.11 & 0.6626 & 1.40 \\
 & RIF & 22.16 & 0.6681 & 1.46 & 22.12 & 0.6692 & 1.41 \\
 & RL & 22.14 & 0.6696 & 1.43 & 22.10 & 0.6484 & 1.40 \\
 & RLTV & 21.89 & 0.6626 & 1.18 & 21.99 & 0.6530 & 1.28 \\
Poisson & Observed & 20.44 & 0.4824 & 0.00 & 20.44 & 0.4824 & -0.00 \\
 & NIF & 8.89 & 0.0606 & -11.55 & 8.89 & 0.0606 & -11.55 \\
 & TRIF & 21.48 & 0.6089 & 1.04 & 21.48 & 0.6089 & 1.04 \\
 & RIF & 21.55 & 0.6210 & 1.11 & 21.53 & 0.6218 & 1.09 \\
 & RL & 21.48 & 0.6235 & 1.04 & 21.21 & 0.5674 & 0.77 \\
 & RLTV & 21.55 & 0.6225 & 1.11 & 21.26 & 0.5692 & 0.82 \\
\bottomrule

\end{tabular}
\end{center}


\clearpage
\subsection{BSDS500 image 100099}
\begin{center}
\includegraphics[width=\linewidth]{atlas-100099.png}
\captionof{figure}{Full-frame and central-ROI comparison for image 100099.}
\end{center}
\begin{center}
\small\setlength{\tabcolsep}{3.3pt}
\captionof{table}{Clipped whole-image results. Gain is relative to the observation (dB).}
\begin{tabular}{llrrrrrr}
\toprule
& & \multicolumn{3}{c}{Python} & \multicolumn{3}{c}{DeconvolutionLab2}\\
\cmidrule(lr){3-5}\cmidrule(lr){6-8}
Noise & Method & PSNR & SSIM & Gain & PSNR & SSIM & Gain\\\midrule
Gaussian & Observed & 29.73 & 0.7259 & 0.00 & 29.73 & 0.7259 & 0.00 \\
 & NIF & 12.62 & 0.0481 & -17.11 & 12.62 & 0.0481 & -17.11 \\
 & TRIF & 30.14 & 0.7374 & 0.41 & 30.14 & 0.7374 & 0.41 \\
 & RIF & 30.13 & 0.7320 & 0.39 & 30.28 & 0.7433 & 0.55 \\
 & RL & 30.12 & 0.7319 & 0.39 & 29.36 & 0.6748 & -0.38 \\
 & RLTV & 31.00 & 0.7985 & 1.27 & 30.04 & 0.7229 & 0.30 \\
Poisson & Observed & 27.98 & 0.5701 & 0.00 & 27.98 & 0.5701 & 0.00 \\
 & NIF & 8.74 & 0.0171 & -19.24 & 8.74 & 0.0171 & -19.24 \\
 & TRIF & 28.26 & 0.6293 & 0.28 & 28.26 & 0.6293 & 0.28 \\
 & RIF & 28.93 & 0.6641 & 0.96 & 29.10 & 0.6764 & 1.13 \\
 & RL & 29.53 & 0.7057 & 1.55 & 27.25 & 0.5345 & -0.72 \\
 & RLTV & 29.98 & 0.7426 & 2.01 & 26.64 & 0.5068 & -1.34 \\
\bottomrule

\end{tabular}
\end{center}


\clearpage
\subsection{BSDS500 image 10081}
\begin{center}
\includegraphics[width=\linewidth]{atlas-10081.png}
\captionof{figure}{Full-frame and central-ROI comparison for image 10081.}
\end{center}
\begin{center}
\small\setlength{\tabcolsep}{3.3pt}
\captionof{table}{Clipped whole-image results. Gain is relative to the observation (dB).}
\begin{tabular}{llrrrrrr}
\toprule
& & \multicolumn{3}{c}{Python} & \multicolumn{3}{c}{DeconvolutionLab2}\\
\cmidrule(lr){3-5}\cmidrule(lr){6-8}
Noise & Method & PSNR & SSIM & Gain & PSNR & SSIM & Gain\\\midrule
Gaussian & Observed & 24.74 & 0.6924 & 0.00 & 24.74 & 0.6924 & 0.00 \\
 & NIF & 12.84 & 0.1167 & -11.89 & 12.84 & 0.1167 & -11.89 \\
 & TRIF & 26.88 & 0.7436 & 2.15 & 26.88 & 0.7436 & 2.15 \\
 & RIF & 27.00 & 0.7389 & 2.26 & 26.98 & 0.7471 & 2.24 \\
 & RL & 26.85 & 0.7331 & 2.11 & 26.59 & 0.6849 & 1.86 \\
 & RLTV & 26.83 & 0.7976 & 2.09 & 26.61 & 0.7220 & 1.88 \\
Poisson & Observed & 23.96 & 0.5285 & 0.00 & 23.96 & 0.5285 & 0.00 \\
 & NIF & 8.44 & 0.0406 & -15.52 & 8.44 & 0.0406 & -15.52 \\
 & TRIF & 25.34 & 0.6183 & 1.38 & 25.34 & 0.6183 & 1.38 \\
 & RIF & 25.69 & 0.6473 & 1.73 & 25.71 & 0.6567 & 1.75 \\
 & RL & 25.68 & 0.6775 & 1.72 & 24.62 & 0.5259 & 0.66 \\
 & RLTV & 26.23 & 0.7458 & 2.27 & 24.46 & 0.5133 & 0.50 \\
\bottomrule

\end{tabular}
\end{center}


\clearpage
\subsection{BSDS500 image 101027}
\begin{center}
\includegraphics[width=\linewidth]{atlas-101027.png}
\captionof{figure}{Full-frame and central-ROI comparison for image 101027.}
\end{center}
\begin{center}
\small\setlength{\tabcolsep}{3.3pt}
\captionof{table}{Clipped whole-image results. Gain is relative to the observation (dB).}
\begin{tabular}{llrrrrrr}
\toprule
& & \multicolumn{3}{c}{Python} & \multicolumn{3}{c}{DeconvolutionLab2}\\
\cmidrule(lr){3-5}\cmidrule(lr){6-8}
Noise & Method & PSNR & SSIM & Gain & PSNR & SSIM & Gain\\\midrule
Gaussian & Observed & 20.15 & 0.5486 & 0.00 & 20.15 & 0.5486 & -0.00 \\
 & NIF & 13.21 & 0.1775 & -6.94 & 13.21 & 0.1775 & -6.94 \\
 & TRIF & 22.11 & 0.6615 & 1.96 & 22.11 & 0.6615 & 1.96 \\
 & RIF & 22.19 & 0.6626 & 2.04 & 22.14 & 0.6652 & 2.00 \\
 & RL & 22.16 & 0.6662 & 2.01 & 22.14 & 0.6422 & 1.99 \\
 & RLTV & 21.87 & 0.6720 & 1.72 & 21.96 & 0.6519 & 1.81 \\
Poisson & Observed & 19.98 & 0.5155 & 0.00 & 19.98 & 0.5155 & -0.00 \\
 & NIF & 9.67 & 0.0806 & -10.31 & 9.67 & 0.0806 & -10.31 \\
 & TRIF & 21.50 & 0.6298 & 1.52 & 21.50 & 0.6298 & 1.52 \\
 & RIF & 21.55 & 0.6390 & 1.57 & 21.52 & 0.6394 & 1.54 \\
 & RL & 21.40 & 0.6416 & 1.42 & 21.30 & 0.5950 & 1.32 \\
 & RLTV & 21.62 & 0.6486 & 1.64 & 21.45 & 0.5948 & 1.47 \\
\bottomrule

\end{tabular}
\end{center}


\clearpage
\subsection{BSDS500 image 101084}
\begin{center}
\includegraphics[width=\linewidth]{atlas-101084.png}
\captionof{figure}{Full-frame and central-ROI comparison for image 101084.}
\end{center}
\begin{center}
\small\setlength{\tabcolsep}{3.3pt}
\captionof{table}{Clipped whole-image results. Gain is relative to the observation (dB).}
\begin{tabular}{llrrrrrr}
\toprule
& & \multicolumn{3}{c}{Python} & \multicolumn{3}{c}{DeconvolutionLab2}\\
\cmidrule(lr){3-5}\cmidrule(lr){6-8}
Noise & Method & PSNR & SSIM & Gain & PSNR & SSIM & Gain\\\midrule
Gaussian & Observed & 20.50 & 0.3371 & 0.00 & 20.50 & 0.3371 & -0.00 \\
 & NIF & 12.92 & 0.1540 & -7.58 & 12.92 & 0.1540 & -7.58 \\
 & TRIF & 21.54 & 0.4660 & 1.04 & 21.54 & 0.4660 & 1.04 \\
 & RIF & 21.56 & 0.4725 & 1.07 & 21.54 & 0.4688 & 1.04 \\
 & RL & 21.59 & 0.4744 & 1.10 & 21.54 & 0.4712 & 1.04 \\
 & RLTV & 21.46 & 0.4473 & 0.97 & 21.47 & 0.4561 & 0.97 \\
Poisson & Observed & 20.34 & 0.3195 & 0.00 & 20.34 & 0.3195 & -0.00 \\
 & NIF & 9.60 & 0.0689 & -10.74 & 9.60 & 0.0689 & -10.74 \\
 & TRIF & 21.17 & 0.4258 & 0.83 & 21.17 & 0.4258 & 0.83 \\
 & RIF & 21.20 & 0.4270 & 0.86 & 21.18 & 0.4239 & 0.84 \\
 & RL & 21.17 & 0.4170 & 0.83 & 21.01 & 0.4085 & 0.67 \\
 & RLTV & 21.30 & 0.4265 & 0.96 & 21.04 & 0.4253 & 0.70 \\
\bottomrule

\end{tabular}
\end{center}


\clearpage
\subsection{BSDS500 image 102062}
\begin{center}
\includegraphics[width=\linewidth]{atlas-102062.png}
\captionof{figure}{Full-frame and central-ROI comparison for image 102062.}
\end{center}
\begin{center}
\small\setlength{\tabcolsep}{3.3pt}
\captionof{table}{Clipped whole-image results. Gain is relative to the observation (dB).}
\begin{tabular}{llrrrrrr}
\toprule
& & \multicolumn{3}{c}{Python} & \multicolumn{3}{c}{DeconvolutionLab2}\\
\cmidrule(lr){3-5}\cmidrule(lr){6-8}
Noise & Method & PSNR & SSIM & Gain & PSNR & SSIM & Gain\\\midrule
Gaussian & Observed & 18.89 & 0.4260 & 0.00 & 18.89 & 0.4260 & -0.00 \\
 & NIF & 12.59 & 0.2252 & -6.29 & 12.59 & 0.2252 & -6.29 \\
 & TRIF & 20.41 & 0.5909 & 1.53 & 20.41 & 0.5909 & 1.53 \\
 & RIF & 20.47 & 0.5983 & 1.58 & 20.42 & 0.5953 & 1.53 \\
 & RL & 20.62 & 0.6095 & 1.73 & 20.63 & 0.6039 & 1.74 \\
 & RLTV & 20.29 & 0.5798 & 1.40 & 20.43 & 0.5894 & 1.55 \\
Poisson & Observed & 18.72 & 0.4018 & 0.00 & 18.72 & 0.4018 & -0.00 \\
 & NIF & 8.87 & 0.0936 & -9.85 & 8.87 & 0.0936 & -9.85 \\
 & TRIF & 19.91 & 0.5426 & 1.19 & 19.91 & 0.5426 & 1.19 \\
 & RIF & 19.92 & 0.5463 & 1.20 & 19.89 & 0.5434 & 1.17 \\
 & RL & 19.91 & 0.5424 & 1.18 & 19.81 & 0.5255 & 1.09 \\
 & RLTV & 20.00 & 0.5500 & 1.28 & 19.97 & 0.5433 & 1.25 \\
\bottomrule

\end{tabular}
\end{center}


\clearpage
\subsection{BSDS500 image 103006}
\begin{center}
\includegraphics[width=\linewidth]{atlas-103006.png}
\captionof{figure}{Full-frame and central-ROI comparison for image 103006.}
\end{center}
\begin{center}
\small\setlength{\tabcolsep}{3.3pt}
\captionof{table}{Clipped whole-image results. Gain is relative to the observation (dB).}
\begin{tabular}{llrrrrrr}
\toprule
& & \multicolumn{3}{c}{Python} & \multicolumn{3}{c}{DeconvolutionLab2}\\
\cmidrule(lr){3-5}\cmidrule(lr){6-8}
Noise & Method & PSNR & SSIM & Gain & PSNR & SSIM & Gain\\\midrule
Gaussian & Observed & 21.58 & 0.6115 & 0.00 & 21.58 & 0.6115 & 0.00 \\
 & NIF & 13.21 & 0.1766 & -8.37 & 13.21 & 0.1766 & -8.37 \\
 & TRIF & 23.47 & 0.7086 & 1.88 & 23.47 & 0.7086 & 1.88 \\
 & RIF & 23.54 & 0.7073 & 1.95 & 23.50 & 0.7105 & 1.92 \\
 & RL & 23.56 & 0.7111 & 1.98 & 23.48 & 0.6836 & 1.90 \\
 & RLTV & 23.39 & 0.7244 & 1.80 & 23.36 & 0.6985 & 1.77 \\
Poisson & Observed & 21.22 & 0.5526 & 0.00 & 21.22 & 0.5526 & 0.00 \\
 & NIF & 8.73 & 0.0605 & -12.50 & 8.73 & 0.0605 & -12.50 \\
 & TRIF & 22.60 & 0.6539 & 1.37 & 22.60 & 0.6539 & 1.37 \\
 & RIF & 22.77 & 0.6671 & 1.54 & 22.76 & 0.6702 & 1.54 \\
 & RL & 22.77 & 0.6828 & 1.55 & 22.32 & 0.6070 & 1.09 \\
 & RLTV & 23.04 & 0.6896 & 1.81 & 22.27 & 0.5992 & 1.04 \\
\bottomrule

\end{tabular}
\end{center}


\clearpage
\subsection{BSDS500 image 103029}
\begin{center}
\includegraphics[width=\linewidth]{atlas-103029.png}
\captionof{figure}{Full-frame and central-ROI comparison for image 103029.}
\end{center}
\begin{center}
\small\setlength{\tabcolsep}{3.3pt}
\captionof{table}{Clipped whole-image results. Gain is relative to the observation (dB).}
\begin{tabular}{llrrrrrr}
\toprule
& & \multicolumn{3}{c}{Python} & \multicolumn{3}{c}{DeconvolutionLab2}\\
\cmidrule(lr){3-5}\cmidrule(lr){6-8}
Noise & Method & PSNR & SSIM & Gain & PSNR & SSIM & Gain\\\midrule
Gaussian & Observed & 28.78 & 0.7856 & 0.00 & 28.78 & 0.7856 & 0.00 \\
 & NIF & 13.06 & 0.0724 & -15.72 & 13.06 & 0.0724 & -15.72 \\
 & TRIF & 30.88 & 0.8155 & 2.10 & 30.88 & 0.8155 & 2.10 \\
 & RIF & 30.90 & 0.8096 & 2.12 & 31.05 & 0.8215 & 2.27 \\
 & RL & 30.96 & 0.8097 & 2.18 & 30.07 & 0.7493 & 1.29 \\
 & RLTV & 31.87 & 0.8774 & 3.09 & 30.74 & 0.7989 & 1.96 \\
Poisson & Observed & 27.67 & 0.6667 & 0.00 & 27.67 & 0.6667 & 0.00 \\
 & NIF & 9.46 & 0.0279 & -18.21 & 9.46 & 0.0279 & -18.21 \\
 & TRIF & 29.19 & 0.7453 & 1.52 & 29.19 & 0.7453 & 1.52 \\
 & RIF & 29.79 & 0.7738 & 2.12 & 29.92 & 0.7842 & 2.25 \\
 & RL & 30.27 & 0.8109 & 2.60 & 28.24 & 0.6562 & 0.57 \\
 & RLTV & 30.92 & 0.8324 & 3.25 & 27.75 & 0.6251 & 0.08 \\
\bottomrule

\end{tabular}
\end{center}


\clearpage
\subsection{BSDS500 image 103078}
\begin{center}
\includegraphics[width=\linewidth]{atlas-103078.png}
\captionof{figure}{Full-frame and central-ROI comparison for image 103078.}
\end{center}
\begin{center}
\small\setlength{\tabcolsep}{3.3pt}
\captionof{table}{Clipped whole-image results. Gain is relative to the observation (dB).}
\begin{tabular}{llrrrrrr}
\toprule
& & \multicolumn{3}{c}{Python} & \multicolumn{3}{c}{DeconvolutionLab2}\\
\cmidrule(lr){3-5}\cmidrule(lr){6-8}
Noise & Method & PSNR & SSIM & Gain & PSNR & SSIM & Gain\\\midrule
Gaussian & Observed & 23.53 & 0.5951 & 0.00 & 23.53 & 0.5951 & -0.00 \\
 & NIF & 13.01 & 0.1275 & -10.52 & 13.01 & 0.1275 & -10.52 \\
 & TRIF & 25.18 & 0.6911 & 1.65 & 25.18 & 0.6911 & 1.65 \\
 & RIF & 25.26 & 0.6952 & 1.73 & 25.26 & 0.6978 & 1.73 \\
 & RL & 25.33 & 0.6990 & 1.80 & 25.15 & 0.6748 & 1.62 \\
 & RLTV & 25.27 & 0.7093 & 1.73 & 25.17 & 0.6867 & 1.63 \\
Poisson & Observed & 23.04 & 0.5121 & 0.00 & 23.04 & 0.5121 & -0.00 \\
 & NIF & 8.68 & 0.0421 & -14.36 & 8.68 & 0.0421 & -14.36 \\
 & TRIF & 24.24 & 0.6103 & 1.20 & 24.24 & 0.6103 & 1.20 \\
 & RIF & 24.53 & 0.6344 & 1.49 & 24.55 & 0.6380 & 1.51 \\
 & RL & 24.61 & 0.6482 & 1.56 & 23.93 & 0.5651 & 0.89 \\
 & RLTV & 24.83 & 0.6672 & 1.79 & 23.85 & 0.5604 & 0.80 \\
\bottomrule

\end{tabular}
\end{center}


\clearpage
\newgeometry{textwidth=6.45in,textheight=9.15in,top=1in,headheight=14pt,headsep=15pt}
\setlength{\headwidth}{\textwidth}
\section{Implementation source code}
\label{app:code}

The excerpts use the original line numbers of the files in the linked GitHub source code folder. Listing~\ref{lst:forward-code} in Section~\ref{sec:implementation} gives the shared forward operator. The listings below follow the simulation, solver, evaluation, and ImageJ reference workflows.

\subsection{PSF construction and noise simulation}

In \file{code/deconv/core.py}, both PSFs are normalized before convolution. The Gaussian observation is floored at zero, and the Poisson observation is scaled back to intensity units after sampling photon counts.

\reportlisting{Python}{deconv/core.py}{33}{44}{
  caption={Unit-sum Gaussian and horizontal motion PSFs.},
  label={lst:psf-code}
}

\reportlisting{Python}{deconv/core.py}{65}{74}{
  caption={Seeded observations with Gaussian or Poisson noise.},
  label={lst:noise-code}
}

\clearpage
\subsection{Deconvolution solvers}

Listings~\ref{lst:solver-helpers}--\ref{lst:iterative-code} show the spectral and TV helpers and the complete \file{solve} function from \file{code/deconv/solvers.py}.

\reportlisting{Python}{deconv/solvers.py}{11}{25}{
  caption={Identity-penalty spectral filter and smoothed TV divergence.},
  label={lst:solver-helpers}
}

\reportlisting{Python}{deconv/solvers.py}{27}{57}{
  caption={Solver entry point, input checks, and FFT setup.},
  label={lst:solver-setup}
}

\clearpage
The direct branch uses the NIF power floor or a Tikhonov penalty. The iterative branch applies the forward and adjoint operators at each step and adds the TV denominator for RLTV.

\reportlisting{Python}{deconv/solvers.py}{58}{73}{
  caption={Direct spectral restoration for NIF, TRIF, and Laplacian RIF.},
  label={lst:spectral-code}
}

\reportlisting{Python}{deconv/solvers.py}{75}{100}{
  caption={Iterative restoration for RL and RLTV, including the ImageJ-matched RL rule.},
  label={lst:iterative-code}
}

\clearpage
\subsection{Parameter selection and test evaluation}

In \file{code/scripts/run_experiments.py}, the validation loop selects the highest mean clipped PSNR for each method and noise model. The script saves the frozen configuration before the test loop scores the held-out images under both policies.

\reportlisting{Python}{scripts/run_experiments.py}{75}{97}{
  caption={Validation-only parameter selection and the Poisson RL diagnostic.},
  label={lst:validation-code}
}

\reportlisting{Python}{scripts/run_experiments.py}{121}{137}{
  caption={Freezing parameters and evaluating the held-out BSDS500 test images.},
  label={lst:test-code}
}

\clearpage
\subsection{Metrics and ImageJ reference}

In \file{code/deconv/metrics.py}, MSE and fixed-peak PSNR use every pixel. The SSIM mean uses the centers of valid \(11\times11\) Gaussian windows. The Java launcher in \file{code/imagej/BatchReference.java} reloads the float TIFF and calls its independent scoring routine after each plugin run.

\reportlisting{Python}{deconv/metrics.py}{6}{18}{
  caption={Python implementation of MSE, PSNR, and local SSIM.},
  label={lst:metrics-code}
}

\reportlisting{Java}{imagej/BatchReference.java}{103}{124}{
  caption={DeconvolutionLab2 invocation, float TIFF decoding, and independent scoring handoff.},
  label={lst:imagej-code}
}

\end{document}
